%% marginfx --- main paper, Annals of Statistics (imsart)
%% Ported from the original article-class source.
\documentclass[aos]{imsart}

\RequirePackage{amsthm,amsmath,amsfonts,amssymb}
\RequirePackage{mathtools}
\RequirePackage[authoryear]{natbib}
\RequirePackage{graphicx}
\RequirePackage{booktabs,multirow,threeparttable,tabularx,rotating}
\RequirePackage{enumerate}
\RequirePackage[colorlinks,citecolor=blue,urlcolor=blue]{hyperref}


\startlocaldefs

\numberwithin{equation}{section}

%% ---- Theorem environments (imsart style) ----
\theoremstyle{plain}
\newtheorem{theorem}{Theorem}
\newtheorem{lemma}{Lemma}
\newtheorem{corollary}{Corollary}
\newtheorem{conjecture}{Conjecture}

\theoremstyle{definition}
\newtheorem{definition}{Definition}
\newtheorem{remark}{Remark}

%% ---- Author macros (carried over unchanged) ----
\newcommand{\norm}[1]{\left\| #1 \right\|}
\newcommand{\R}{\mathbb{R}}
\newcommand{\E}{\mathbb{E}}
\newcommand{\Lp}{L^p(\Omega)}
\newcommand{\Lone}{L^1(\Omega)}
\newcommand{\Wonep}{W^{1,p}(\Omega)}
\newcommand{\Wtwop}{W^{2,p}(\Omega)}
\newcommand{\BV}{BV(\Omega)}
\newcommand{\cM}{\mathcal{M}(\Omega)^d}
\newcommand{\fn}{\hat{f}_n}
\newcommand{\xrightarrowp}{\xrightarrow{p}}
\newcommand{\wstar}{\xrightarrow{w*}}
\newcommand{\Prob}{\mathbb{P}}
\newcommand{\abs}[1]{\left| #1 \right|}

\endlocaldefs

\begin{document}

\begin{frontmatter}

\title{Model-Agnostic Interpretation of Machine Learning\\via Average Marginal Effects}
\runtitle{Model-Agnostic Interpretation via Average Marginal Effects}

\begin{aug}
\author[A]{\fnms{Jason}~\snm{Parker}\ead[label=e1]{jparker588@gmail.com}}

\address[A]{Artificial Intelligence and Data Transformation, Southwest Airlines\printead[presep={,\ }]{e1}}
\end{aug}

\begin{abstract}
We propose a framework for interpreting any supervised machine learning model
through average marginal effects, the standard inferential tool in the
econometrics of nonlinear models. The procedure is agnostic to the choice of
learner and returns a coefficient table directly analogous to parametric
regression output, with standard errors obtained by a nonparametric bootstrap
that refits the model on each replicate; the classical generalized linear model
result is nested as a special case. Our theoretical contribution is twofold. We
establish gradient consistency for neural networks in the Sobolev space
$W^{1,p}$ and for gradient boosted models and Mondrian forests in the space of
functions of bounded variation, and identify the obstruction to extending the
result to Breiman forests. We then prove bootstrap validity for all three
classes via Hadamard differentiability of the average marginal effect functional
and the functional delta method, giving the first rigorous inferential
foundation for machine learning average marginal effects. The
interpretability--accuracy tradeoff is therefore not fundamental but an artifact
of requiring parametric assumptions to make interpretation rigorous. Experiments
on three datasets confirm that estimates are stable across learners, with
Breiman forests the exception, as the theory predicts.
\end{abstract}

\begin{keyword}[class=MSC]
\kwdgroup[type=primary]{\kwd{62G08}
\kwd{62G09}}
\kwdgroup[type=secondary]{\kwd{62G20}}
\end{keyword}

\begin{keyword}
\kwd{average marginal effects}
\kwd{model-agnostic interpretability}
\kwd{bootstrap validity}
\kwd{Hadamard differentiability}
\kwd{Mondrian forests}
\end{keyword}

\end{frontmatter}

\tableofcontents

\section{Introduction}

The rise of flexible machine learning methods has created a persistent tension in
empirical research. Models such as random forests \citep{breiman2001rf}, gradient
boosted trees \citep{chen2016xgboost}, and deep neural networks \citep{hastie2009esl}
routinely achieve predictive accuracy that parametric models cannot match. Yet the
standard tools of statistical inference --- coefficient estimates, standard errors,
hypothesis tests --- apply naturally only to parametric models whose functional form
is specified in advance. \citet{breiman2001cultures} framed this as a conflict between
two cultures of statistical modeling: one oriented toward interpretable
data-generating models, the other toward accurate algorithmic prediction. The
conventional resolution has been to treat accuracy and interpretability as opposing
ends of a spectrum, accepting reduced interpretability as the price of improved
prediction. This paper argues that the tradeoff is not fundamental. It is an artifact
of requiring parametric assumptions to make interpretation rigorous, and it disappears
once that requirement is dropped.

The machine learning literature has produced a rich set of post-hoc interpretability
methods that attempt to explain black-box model predictions without parametric
assumptions. Partial dependence plots \citep{friedman2001pdp} visualize the marginal
relationship between a feature and the predicted outcome by averaging over the
empirical distribution of other features. LIME \citep{ribeiro2016lime} fits a local
linear approximation to the prediction function in the neighborhood of individual
observations. SHAP \citep{lundberg2017shap} assigns each feature a contribution to
each prediction based on Shapley values from cooperative game theory, providing a
unified framework that nests several earlier methods. These tools are widely used in
practice and are surveyed comprehensively in \citet{molnar2025} and
\citet{doshivelez2017interpretability}. However, they share a fundamental limitation:
they produce point estimates with no associated sampling distributions. A SHAP value,
a PDP slope, and a LIME coefficient are descriptive quantities --- there is no
standard error, no confidence interval, and no basis for a hypothesis test. A
researcher cannot determine whether a SHAP value is statistically distinguishable
from zero, or whether two features have significantly different effects, or whether
the effect of a variable differs across model specifications. For applied researchers
who need to draw defensible conclusions from their models, this is a serious
constraint.

Two prior contributions are closest in spirit to the present paper.
\citet{scholbeck2024fme} propose forward marginal effects for nonlinear prediction
functions as a model-agnostic interpretability method, establishing a general
framework for computing feature effects for any supervised learner. Their framework
is descriptive rather than inferential: it characterizes how feature effects vary
across the covariate space but does not provide standard errors, confidence intervals,
or hypothesis tests, and it explicitly argues against summarizing feature effects in
a single metric such as the average marginal effect. The present paper takes the
opposite position: we argue that the average marginal effect is precisely the right
summary quantity, that it has a rigorous population-level interpretation inherited
from the econometric literature, and that valid inference is achievable via a
nonparametric bootstrap with full model refitting. \cite{lowe2024fmeffects}
provide the \texttt{fmeffects} R package implementing this descriptive framework,
demonstrating the practical demand for marginal-effect-style interpretation of
machine learning models, but without the theoretical justification --- no gradient
consistency, no bootstrap validity, no characterization of when the marginal effects
are well-defined --- that the present paper provides through the Hadamard
differentiability of the AME functional and the functional delta method.

The econometric literature has solved exactly this inferential problem for parametric
models through average marginal effects \citep{greene2012, wooldridge2010}. For a
generalized linear model with inverse link function $\mu$ and linear predictor
$\eta = \mathbf{x}^\top \boldsymbol{\beta}$, the average marginal effect of variable
$j$ is $\mathrm{AME}_j = \beta_j \cdot \mathbb{E}[\mu'(\eta)]$ --- the regression
coefficient scaled by the average derivative of the inverse link, averaged over the
empirical distribution of the linear predictor. For the identity link this reduces
to the OLS coefficient. For the logistic link it averages the logistic density over
the sample. The result is a quantity measured in the natural units of the outcome,
interpretable as the average change in the predicted outcome associated with a
one-unit change in $x_j$, with standard errors obtained by the delta method or
bootstrap. This is the output of \texttt{margins} in R and Stata \citep{leeper2021},
and it is the standard form in which marginal effects are reported in applied
economics, political science, sociology, and public health. The limitation is that
it has been restricted to parametric models: to compute $\partial f / \partial x_j$
analytically, one must know the functional form of $f$.

This paper removes that restriction. We propose a framework for computing average
marginal effects for any supervised machine learning model, requiring only that the
prediction function $\hat{f}$ can be evaluated at a point. The average marginal
effect of variable $j$ is defined as the expectation of the partial derivative of
$f$ with respect to $x_j$, averaged over the empirical distribution of the features.
For models where the derivative is not available in closed form --- which includes
tree-based models, for which classical derivatives do not exist everywhere --- we
estimate it via a centered finite difference with an adaptive step size scaled to
the spread of each feature. Standard errors are obtained via a nonparametric
bootstrap \citep{efron1979, efrontibshirani1993} that refits the model, including
hyperparameter selection, on each bootstrap sample, capturing uncertainty about model
specification as well as estimation. The result is a coefficient table --- estimate,
standard error, and significance level --- directly analogous to the output of a
parametric regression but with no restrictions on the functional form of $\hat{f}$.
The framework nests the classical GLM result as a special case: when $\hat{f}$ is a
logistic regression, the estimator recovers the classical AME automatically. When
$\hat{f}$ is XGBoost or a random forest or a neural network, the same formula applies
and produces output on the same scale and with the same interpretation.

The theoretical contribution of the paper proceeds in two parts. The first establishes
gradient consistency for three classes of supervised learners, providing the
asymptotic foundation for the proposed estimator. For neural networks, which live in
Sobolev space $W^{1,p}$, gradient consistency follows from uniform $W^{2,p}$
boundedness and $L^p$ consistency via the Rellich--Kondrachov compact embedding
theorem \citep{evans2010pde, hornik1990universal}. For gradient boosted models with
tree base learners, the natural function space is the space of functions of bounded
variation \citep{ambrosio2000bv}, where gradients exist as Radon measures and
consistency is established via the BV compactness theorem \citep{buhlmann2003boosting,
	mason2000boosting}. For random forests, gradient consistency is proved for the
Mondrian forest variant \citep{lakshminarayanan2014mondrian, mourtada2017mondrian}
and conjectured for Breiman forests \citep{breiman2001rf, scornet2015forests}, with
the obstruction to a full proof identified precisely. The second part establishes
bootstrap validity for each of the three estimator classes. The key result is the
Hadamard differentiability of the AME functional
\citep[Chapter~3.9]{vandervaart1996weak}, which shows that the AME varies smoothly
under perturbations of the underlying estimator in the appropriate function space
sense. Combined with rate conditions for each estimator class and the functional delta
method, this yields bootstrap validity as a corollary for neural networks, gradient
boosted models, and Mondrian forests. Bootstrap validity for Breiman forests remains
open for the same reason that gradient consistency is harder to establish --- the
data-dependent splitting rule couples partition geometry to response variation in a
way that resists clean analysis.

We validate the framework empirically on three datasets spanning binary classification
in social science, binary classification in finance, and continuous regression: the
UCI Adult Income dataset \citep{kohavi1996, UCI}, the UCI Credit Card Default dataset
\citep{yeh2009, UCI}, and the Ames Housing dataset \citep{decock2011}. Across all
three applications we compare AMEs from logistic regression, random forests
\citep{breiman2001rf}, XGBoost \citep{chen2016xgboost}, and neural networks against
SHAP values \citep{lundberg2017shap} and partial dependence plot slopes
\citep{friedman2001pdp}. The AME and PDP estimates are in close agreement throughout,
while SHAP values are on a different scale and in several cases have the wrong sign
--- a systematic divergence that is most visible in the regression setting where units
are dollars and the magnitudes are large enough to make the discrepancy concrete.
Because AMEs come with standard errors and SHAP values do not, the AME framework
supports conclusions that the SHAP framework cannot. An additional empirical finding
is that Breiman random forests produce less stable AME estimates than the other
estimator classes, consistent with the theoretical result that gradient consistency
for Breiman forests is harder to establish than for Mondrian forests. Practitioners
who require random forest AMEs are therefore advised to consider Mondrian forests,
for which both gradient consistency and bootstrap validity are proved.

The remainder of the paper is organized as follows. Section~2 defines the estimand, presents the estimator, and establishes its relationship to classical average marginal effects for the GLM family. Section~3 establishes gradient consistency and bootstrap validity for neural networks, gradient boosted models, and Mondrian forests, and identifies the obstruction to extending these results to Breiman forests. Section~4 presents the three empirical applications. Section~5 presents simulation evidence on 	finite-sample performance. Section~6 concludes.

\section{Method}

We propose a framework for interpreting any supervised machine learning model through the lens of average marginal effects. The framework is agnostic to the choice of learner --- the same estimation procedure applies whether $\hat{f}$ is a neural network, a gradient boosted model, a random forest, or any other supervised learner. Crucially, the framework produces standard errors via a nonparametric bootstrap, enabling hypothesis tests and confidence intervals that existing model-agnostic interpretability methods do not support. SHAP values and permutation-based feature importance measures provide point estimates of variable importance but no basis for inference. The average marginal effect estimator proposed here yields a coefficient table --- estimate, standard error, and significance --- directly analogous to the output of a parametric regression, with no restrictions on the functional form of $\hat{f}$.

\subsection{Estimand}

Let $\mathbf{x}_i = (x_{i1}, \dots, x_{ip})^\top \in \mathbb{R}^p$ denote the feature vector for observation $i$, and let $y_i \in \mathbb{R}$ denote the outcome. We observe $n$ independent draws $\{(\mathbf{x}_i, y_i)\}_{i=1}^n$ from some population distribution $P$. Let $f: \mathbb{R}^p \to \mathbb{R}$ denote the conditional expectation function $f(\mathbf{x}) = E[Y \mid \mathbf{x}]$. For binary outcomes, $f(\mathbf{x}) = P(Y = 1 \mid \mathbf{x})$ is the conditional probability of the positive class; for continuous outcomes, $f(\mathbf{x}) = E[Y \mid \mathbf{x}]$ directly. In both cases $f$ is the same object --- the conditional expectation of $Y$ given $\mathbf{x}$ --- and the method requires no modification across outcome types.

The marginal effect of variable $j$ at observation $i$ is the partial derivative of $f$ evaluated at $\mathbf{x}_i$:
\begin{equation}
	m_{ij} = \frac{\partial f}{\partial x_j}\bigg|_{\mathbf{x}_i}
\end{equation}
The population average marginal effect of variable $j$ is:
\begin{equation}
	\mathrm{AME}_j = \mathbb{E}\left[\frac{\partial f}{\partial x_j}(\mathbf{x})\right]
\end{equation}
where the expectation is taken over the marginal distribution of $\mathbf{x}$ under $P$. This is the estimand. The conditions under which $\partial f / \partial x_j$ exists and under which the sample estimator defined below converges to $\mathrm{AME}_j$ are established in Section~4.

\subsection{Estimation}

Given an estimator $\hat{f}$ of $f$, the sample analog is:
\begin{equation}
	\widehat{\mathrm{AME}}_j = \frac{1}{n} \sum_{i=1}^n \frac{\partial \hat{f}}{\partial x_j}\bigg|_{\mathbf{x}_i}
\end{equation}
For most supervised learners the partial derivative $\partial \hat{f} / \partial x_j$ is not available in closed form, and for tree-based models it does not exist everywhere in the classical sense. We therefore estimate it via a centered finite difference:
\begin{equation}
	\frac{\partial \hat{f}}{\partial x_j}\bigg|_{\mathbf{x}_i} \approx \frac{\hat{f}\!\left(\mathbf{x}_i^{(j,\, x_{ij} + h_j)}\right) - \hat{f}\!\left(\mathbf{x}_i^{(j,\, x_{ij} - h_j)}\right)}{2h_j}
\end{equation}
where $\mathbf{x}_i^{(j,v)}$ denotes the vector $\mathbf{x}_i$ with the $j$-th component replaced by $v$. The centered finite difference is preferred over a one-sided difference because its approximation error is $O(h_j^2)$ rather than $O(h_j)$, which becomes material when $h_j$ is not negligibly small.

The step size $h_j$ is chosen adaptively as:
\begin{equation}
	h_j = \max\!\left(10^{-4},\; 0.05 \cdot \hat{\sigma}_j\right)
\end{equation}
where $\hat{\sigma}_j$ is the sample standard deviation of the $j$-th feature. Scaling $h_j$ to the spread of $x_j$ ensures the finite difference is neither so small as to introduce numerical instability nor so large as to obscure local behavior. The floor of $10^{-4}$ guards against degenerate cases where $\hat{\sigma}_j \approx 0$.

For binary or categorical variables the partial derivative does not exist and the finite difference is replaced by a contrast. The marginal effect for observation $i$ is:
\begin{equation}
	m_{ij} = \hat{f}\!\left(\mathbf{x}_i^{(j,1)}\right) - \hat{f}\!\left(\mathbf{x}_i^{(j,0)}\right)
\end{equation}
which measures the change in predicted outcome when $x_j$ is switched from 0 to 1 holding all other variables at their observed values. The sample average marginal effect is then $\widehat{\mathrm{AME}}_j = n^{-1}\sum_{i=1}^n m_{ij}$, with interpretation identical to the continuous case: a change in $x_j$ from 0 to 1 is associated with a $\widehat{\mathrm{AME}}_j$ change in the predicted outcome, on average across the sample.

For binary classification models, $\hat{f}$ must return the estimated conditional probability $\hat{P}(Y=1 \mid \mathbf{x})$ rather than the predicted class label. Class labels are discrete and their finite differences are identically zero; the method requires a continuous output. This is the predicted probability of the positive class, available as \texttt{predict\_proba(X)[:,1]} in scikit-learn and equivalent interfaces in other frameworks. The AME is then interpreted as the change in predicted probability associated with a one-unit change in $x_j$, on average across the sample.

\subsection{Relationship to Classical Average Marginal Effects}

Average marginal effects are a standard tool in the econometrics of nonlinear models \citep{greene2012, wooldridge2010}. For parametric models the estimator $\widehat{\mathrm{AME}}_j$ defined above reduces exactly to the classical result. Consider a generalized linear model with inverse link function $\mu$ and linear predictor $\eta_i = \mathbf{x}_i^\top \boldsymbol{\beta}$, so that $f(\mathbf{x}_i) = \mu(\eta_i)$. The marginal effect of variable $j$ at observation $i$ is:
\begin{equation}
	\frac{\partial f}{\partial x_j}\bigg|_{\mathbf{x}_i} = \beta_j \cdot \mu'(\eta_i)
\end{equation}
and the population average marginal effect is $\mathrm{AME}_j = \beta_j \cdot E[\mu'(\eta)]$. For the identity link $\mu'(\eta) = 1$ and $\mathrm{AME}_j = \beta_j$ --- the marginal effect is constant across observations and equals the OLS coefficient. For the logistic link $\mu'(\eta) = \Lambda(\eta)(1 - \Lambda(\eta))$ and the AME averages the logistic density over the sample \citep{greene2012}. For the probit link $\mu'(\eta) = \phi(\eta)$, the standard normal density evaluated at the linear predictor \citep{wooldridge2010}.

In each case $\widehat{\mathrm{AME}}_j$ recovers the classical result automatically --- the finite difference applied to the model output numerically approximates $\beta_j \cdot \mu'(\eta_i)$ without requiring knowledge of the link function. The framework therefore nests the entire GLM family as a special case and reduces to familiar output when the model is parametric. This backward compatibility is what the \texttt{marginaleffects} package in R and Python implements \citep{arelbundock2024marginaleffects}; the present framework extends that implementation with formal asymptotic guarantees for machine learning models.

\subsection{Inference}

Standard errors are obtained via a nonparametric bootstrap \citep{efron1979, efrontibshirani1993}. Let $\mathcal{D} = \{(\mathbf{x}_i, y_i)\}_{i=1}^n$ denote the observed sample and let $\hat{\lambda}$ denote the hyperparameters of $\hat{f}$ selected by cross-validation on $\mathcal{D}$. For $b = 1, \dots, B$:
\begin{enumerate}[(1)]
	\item Draw a bootstrap sample $\mathcal{D}^{(b)}$ by sampling $n$ observations from $\mathcal{D}$ with replacement, sampling $(\mathbf{x}_i, y_i)$ pairs jointly to preserve the joint distribution of features and outcome.
	\item Initialize the hyperparameter search at $\hat{\lambda}$ and refit $\hat{f}^{(b)}$ on $\mathcal{D}^{(b)}$ with hyperparameters $\hat{\lambda}^{(b)}$ selected by cross-validation.
	\item Compute $\widehat{\mathrm{AME}}_j^{(b)}$ using $\hat{f}^{(b)}$ evaluated at the original sample $\mathcal{D}$.
\end{enumerate}
The bootstrap standard error is:
\begin{equation}
	\widehat{\mathrm{se}}_j = \sqrt{\frac{1}{B-1}\sum_{b=1}^B \left(\widehat{\mathrm{AME}}_j^{(b)} - \bar{\mathrm{AME}}_j\right)^2}
\end{equation}
where $\bar{\mathrm{AME}}_j = B^{-1}\sum_{b=1}^B \widehat{\mathrm{AME}}_j^{(b)}$.

Several design choices warrant comment. Sampling $(\mathbf{x}_i, y_i)$ pairs jointly preserves the joint distribution of features and outcome that $\hat{f}$ is trained to approximate; resampling features and outcomes independently would destroy this dependence structure. Refitting $\hat{f}^{(b)}$ with hyperparameter selection on each bootstrap sample ensures that the bootstrap standard errors capture uncertainty about model specification, not only estimation uncertainty conditional on a fixed architecture. Initializing the hyperparameter search at $\hat{\lambda}$ rather than from scratch makes this computationally feasible --- the original model serves as a warm start, substantially reducing the cost of cross-validation on each replicate. Finally, evaluating $\hat{f}^{(b)}$ at the original sample $\mathcal{D}$ rather than the bootstrap sample $\mathcal{D}^{(b)}$ follows standard practice for smooth functionals of the empirical distribution \citep{efrontibshirani1993} and ensures that $\widehat{\mathrm{AME}}_j^{(b)}$ depends on the bootstrap draw only through $\hat{f}^{(b)}$, which simplifies the asymptotic analysis in Section~4. Analysis of the variant that evaluates at $\mathcal{D}^{(b)}$ is left to future work.

\section{Asymptotic Theory}
\label{sec:theory}

\subsection{Overview}

The method agnostic interpretation framework rests on a theoretical foundation that
connects the gradients of machine learning estimators to the gradients of the true
underlying function. This section establishes that foundation. We show that for three
estimator classes---neural networks, gradient boosted models, and random
forests---the gradient of the estimator $\fn$ converges asymptotically to the gradient
of the true function $f$, and that bootstrap standard errors for the resulting average
marginal effects are asymptotically valid. The average marginal effects presented in
Section~3 are therefore consistently estimating population quantities with well defined
interpretations, not heuristic approximations, and inference on those quantities is
justified.

The connection between the theory and the method is direct. The average marginal
effect of variable $j$ is defined as:
\begin{equation}
	\widehat{AME}_j = \frac{1}{n} \sum_{i=1}^n \frac{\partial \fn}{\partial x_j}(x_i)
	\label{eq:ame}
\end{equation}
For this quantity to be meaningful as an estimate of a population parameter, two
things must be true. First, the gradient $\frac{\partial \fn}{\partial x_j}$ must
exist and be well defined---it must be estimating something real rather than an
artifact of the particular model fitted. Second, it must converge to the
corresponding gradient of the true function $f$ as the sample size grows. Without
these guarantees, the average marginal effects are descriptive statistics about the
fitted model rather than estimates of population quantities. The theorems in this
section provide exactly these guarantees. The corollaries then establish that
bootstrap standard errors correctly characterize the uncertainty in those estimates.

The key insight is that the appropriate notion of gradient consistency depends on the
function space the estimator naturally inhabits, which in turn reflects the inductive
bias of the estimator class. This is not merely a technical observation---it reflects
something fundamental about what each estimator class is doing. Neural networks with
smooth activation functions produce smooth estimators whose gradients are ordinary
$L^p$ functions, and the appropriate framework is Sobolev space theory
\citep{evans2010pde}. Random forests produce piecewise constant estimators whose
gradients are concentrated on leaf boundaries rather than distributed across the
domain, and the appropriate framework is the theory of functions of bounded variation
\citep{ambrosio2000bv}. Gradient boosted models occupy a middle ground depending on
the base learner. The average marginal effects computed in Section~3 are the same
operation---averaging the gradient over the sample---but the theoretical justification
for that operation differs across estimator classes in a way that reflects their
fundamentally different structures.

This observation also clarifies why existing approaches to model interpretation lack
the theoretical foundation developed here. Methods such as SHAP values and partial
dependence plots are not computing gradients of the estimator in any well defined
function space sense, and consequently there is no analog of the results below that
would establish what population quantity they are estimating or at what rate. The
average marginal effect framework, by contrast, is precisely defined as a gradient
operation, and the results below establish that this operation is doing what
practitioners intuitively expect it to do.

Full proofs of all results are provided in the Appendix.

\subsection{Notation and Preliminaries}

Let $\Omega \subset \R^d$ be a bounded Lipschitz domain and let $f: \Omega \to \R$
be the true underlying function relating covariates to outcomes. Let
$\fn: \Omega \to \R$ denote an estimator of $f$ constructed from $n$ observations
$\{(x_i, y_i)\}_{i=1}^n$. We assume throughout that $\Omega$ is bounded and
Lipschitz, which ensures the compactness and embedding results we rely on are
available. These are mild regularity conditions satisfied by any compact covariate
space with reasonable boundary behavior.

We work with two function spaces depending on the estimator class. The choice of
function space is not arbitrary---each space is the natural habitat of the estimator
class in question, in the sense that estimators of that class belong to the space
with probability one under mild conditions.

The \textbf{Sobolev space} $W^{k,p}(\Omega)$ for $k \geq 1$ and $p \geq 1$
consists of functions whose weak derivatives up to order $k$ are in $\Lp$, with
norm:
\begin{equation}
	\norm{g}_{W^{k,p}} = \sum_{|\alpha| \leq k} \norm{\partial^\alpha g}_{L^p}
\end{equation}
The gradient $\nabla g \in L^p(\Omega)^d$ is a vector valued $L^p$ function for any
$g \in \Wonep$. Sobolev spaces are the natural setting for smooth estimators because
they quantify not just how well the function is approximated in $L^p$ but how well
its derivatives are approximated simultaneously. This is exactly the property needed
to justify gradient based interpretation---$L^p$ consistency of $\fn$ alone is
insufficient because differentiation is an unbounded operator on $L^p$, and
convergence of a function does not in general imply convergence of its derivatives
without additional regularity \citep[Chapter~5]{evans2010pde}.

The \textbf{space of functions of bounded variation} $\BV$ consists of functions
$g \in \Lone$ whose distributional gradient $Dg$ is a vector valued Radon measure
with finite total variation $|Dg|(\Omega) < \infty$. The $BV$ norm is:
\begin{equation}
	\norm{g}_{BV} = \norm{g}_{L^1} + |Dg|(\Omega)
\end{equation}
The $BV$ framework extends the notion of gradient to functions that are not
differentiable in the classical or even weak sense---precisely the situation with
tree based estimators, which have jump discontinuities at leaf boundaries. For
piecewise constant functions the distributional gradient is a sum of surface measures
concentrated on the boundaries between constant regions, capturing the intuition that
the ``derivative'' of a step function is a spike at the jump location. This makes
$BV$ the natural space for random forests and boosting with tree base learners. The
total variation $|Dg|(\Omega)$ measures the total size of all jumps, weighted by the
boundary surface area---a quantity that is finite and well controlled for any
piecewise constant function with finitely many pieces. The standard reference for
the theory of $BV$ functions is \citet{ambrosio2000bv}.

The relationship between the two spaces is important. If $f \in \Wonep$ for
$p \geq 1$, then $f \in \BV$ with $Df = \nabla f \cdot \mathcal{L}^d$---that is,
the distributional gradient of $f$ is absolutely continuous with respect to Lebesgue
measure with density $\nabla f$. This embedding $W^{1,p} \hookrightarrow BV$ means
that assuming $f \in \Wonep$ is sufficient for both the Sobolev and BV results below.
We make this assumption throughout.

\begin{definition}[Well-Behaved Estimator, Sobolev Sense]
	\label{def:wellbehaved_sobolev}
	We say that $\fn$ is \emph{well-behaved in the Sobolev sense} if:
	\begin{equation}
		\sup_n \norm{\fn}_{W^{2,p}} < \infty \quad \text{and} \quad
		\norm{\fn - f}_{L^p} \xrightarrowp 0
	\end{equation}
\end{definition}

\begin{definition}[Well-Behaved Estimator, BV Sense]
	\label{def:wellbehaved_bv}
	We say that $\fn$ is \emph{well-behaved in the BV sense} if:
	\begin{equation}
		\sup_n \E\left[|D\fn|(\Omega)\right] < \infty \quad \text{and} \quad
		\norm{\fn - f}_{L^1} \xrightarrowp 0
	\end{equation}
\end{definition}

In both cases the well-behaved condition has the same structure: a uniform
boundedness condition in a strong space combined with consistency in a weaker space.
The uniform boundedness condition is the key regularity requirement---it says that
the estimator sequence does not become increasingly irregular as $n$ grows, even as
it fits the data more closely. This condition is verifiable from properties of the
estimator class and is established for each class in the theorems below. The
consistency condition is the standard requirement that the estimator converges to the
truth and is available from existing results in each case.

We also state the rate condition required for the bootstrap corollaries. Let $P$
denote the marginal distribution of the covariates $x$.

\begin{definition}[Rate Condition]
	\label{def:rate}
	We say that $\fn$ satisfies the \emph{rate condition} if:
	\begin{equation}
		\sqrt{n}\left(\fn - f\right) = O_p(1)
	\end{equation}
	in the appropriate norm --- $\Wonep$ for estimators well-behaved in the Sobolev
	sense, and $\BV$ for estimators well-behaved in the BV sense.
\end{definition}

The rate condition requires that $\fn$ converges to $f$ at the parametric rate
$n^{-1/2}$. This is stronger than consistency alone and is the condition that makes
the functional delta method applicable. It is verified for each estimator class in
the corollaries below.

\subsection{Core Lemmas}

The theoretical results rest on five lemmas. Lemmas~\ref{lem:sobolev_existence}
through~\ref{lem:bv_consistency} establish derivative existence and consistency in
each function space. Lemma~\ref{lem:hadamard} establishes the Hadamard
differentiability of the AME functional, which is the key condition for bootstrap
validity. All five lemmas are proved in the Appendix and applied across all three
estimator classes.

The parallel structure of Lemmas~\ref{lem:sobolev_existence}
through~\ref{lem:bv_consistency} reflects the parallel structure of the Sobolev
and BV frameworks and is the formal expression of the analogy between smooth and
piecewise constant estimators. Lemma~\ref{lem:hadamard} then sits above both
frameworks as the unifying result that connects gradient consistency to bootstrap
validity.

\begin{lemma}[Derivative Existence in Sobolev Spaces]
	\label{lem:sobolev_existence}
	Let $\fn \in \Wtwop$ for some $p > 1$. Then $\nabla \fn \in W^{1,p}(\Omega)^d$
	exists as a well defined $L^p$ function and satisfies:
	\begin{equation}
		\norm{\nabla \fn}_{L^p} \leq \norm{\fn}_{W^{1,p}} \leq \norm{\fn}_{W^{2,p}}
	\end{equation}
\end{lemma}

Lemma~\ref{lem:sobolev_existence} is essentially definitional---membership in
$\Wtwop$ implies by definition that the first and second weak derivatives exist and
are in $L^p$. The content of the lemma is the norm bound, which establishes that
the gradient operator $\nabla: W^{2,p} \to W^{1,p}$ is bounded and linear. This
boundedness is what makes the subsequent consistency argument work---a bounded linear
operator preserves convergence, so once we have convergence of $\fn$ in the right
space, convergence of $\nabla \fn$ follows automatically. The reason we require
$\fn \in \Wtwop$ rather than just $\Wonep$ is that we need one additional degree of
differentiability to apply the compactness argument in
Lemma~\ref{lem:sobolev_consistency}. The relevant theory of Sobolev spaces and their
norms is developed in \citet[Chapter~5]{evans2010pde}.

\begin{lemma}[Derivative Consistency in Sobolev Spaces]
	\label{lem:sobolev_consistency}
	Let $f \in \Wonep$ and suppose $\fn$ is well-behaved in the Sobolev sense. Then:
	\begin{equation}
		\norm{\nabla \fn - \nabla f}_{L^p} \xrightarrowp 0
	\end{equation}
\end{lemma}

Lemma~\ref{lem:sobolev_consistency} is the harder result and the one that does the
real work in the Sobolev setting. The difficulty is that $L^p$ consistency of $\fn$
alone does not imply $L^p$ consistency of $\nabla \fn$---differentiation is an
unbounded operator on $L^p$ and limits and derivatives cannot be exchanged without
additional structure. The uniform $\Wtwop$ boundedness condition provides that
structure. The proof applies the Rellich--Kondrachov compact embedding theorem, which
states that a uniformly bounded sequence in $\Wtwop$ has a subsequence converging
strongly in $\Wonep$ \citep[Theorem~5,~Section~5.7]{evans2010pde}. Combined with
the $L^p$ consistency condition identifying the limit as $f$, this gives strong
convergence of $\fn$ to $f$ in $\Wonep$, from which $L^p$ convergence of the
gradients follows immediately from the boundedness of $\nabla$ established in
Lemma~\ref{lem:sobolev_existence}. The key insight is that the uniform boundedness
condition ensures the estimator sequence lives in a compact subset of $\Wonep$,
preventing the gradients from escaping to infinity even as the estimator fits the
data more closely.

\begin{lemma}[Derivative Existence in BV]
	\label{lem:bv_existence}
	Let $\fn \in \BV$. Then the distributional gradient $D\fn$ exists as a well defined
	vector valued Radon measure in $\cM$ with:
	\begin{equation}
		|D\fn|(\Omega) < \infty
	\end{equation}
	For piecewise constant estimators with leaf partition $\{R_j\}$ and leaf means
	$\{\bar{y}_j\}$, the distributional gradient takes the explicit form:
	\begin{equation}
		D\fn = \sum_{j \sim k} (\bar{y}_k - \bar{y}_j) \cdot \nu_{jk} \cdot
		\mathcal{H}^{d-1}\lfloor_{\partial R_j \cap \partial R_k}
	\end{equation}
	where $\nu_{jk}$ is the unit normal between adjacent leaves and $\mathcal{H}^{d-1}$
	is the $(d-1)$-dimensional Hausdorff measure on leaf boundaries.
\end{lemma}

Lemma~\ref{lem:bv_existence} is the BV analog of Lemma~\ref{lem:sobolev_existence}
and is similarly structural in nature. The first part---that $BV$ membership implies
the distributional gradient exists as a Radon measure---follows from the definition
of $BV$ as developed in \citet[Chapter~3]{ambrosio2000bv}. The content of the lemma
is the explicit formula for piecewise constant estimators, which makes concrete what
the distributional gradient actually looks like for a tree based model. The formula
says that the gradient is not a function at all but a measure concentrated on the
boundaries between leaves, with mass proportional to the jump in response values
across the boundary weighted by the boundary surface area. This is the precise sense
in which a random forest or boosted tree has a gradient---not a pointwise derivative,
which does not exist, but a distributional object that captures where and how much
the estimator changes. This explicit representation is also what makes the total
variation $|D\fn|(\Omega)$ computable in terms of observable quantities---the leaf
means and the partition geometry---which is essential for verifying the well-behaved
condition in the theorems below.

\begin{lemma}[Derivative Consistency in BV]
	\label{lem:bv_consistency}
	Let $f \in \Wonep$ for some $p > 1$, so that $f \in \BV$ with
	$Df = \nabla f \cdot \mathcal{L}^d$. Suppose $\fn$ is well-behaved in the BV sense.
	Then:
	\begin{equation}
		D\fn \wstar Df \quad \text{in } \cM
	\end{equation}
	If additionally $\E[|D\fn|(\Omega)] \to |Df|(\Omega)$, then convergence is strict.
\end{lemma}

Lemma~\ref{lem:bv_consistency} is the hardest of the four function space lemmas and
the BV analog of Lemma~\ref{lem:sobolev_consistency}. The convergence is necessarily
weaker than in the Sobolev case---weak-$*$ convergence of measures rather than strong
$L^p$ convergence---because the $BV$ space is not reflexive and the stronger norm
convergence is not available in general. Weak-$*$ convergence means that for any
continuous test function $\phi$:
\begin{equation}
	\int_\Omega \phi \, d(D\fn) \to \int_\Omega \phi \, d(Df)
\end{equation}
which is the appropriate notion of convergence for the distributional gradients of
piecewise constant estimators. In the context of average marginal effects, this means
that integrated gradient information---which is exactly what the AME computes---
converges to the correct population quantity. The proof applies the BV compactness
theorem of \citet[Theorem~3.23]{ambrosio2000bv}, which is the analog of
Rellich--Kondrachov in the BV setting: a uniformly bounded sequence in $BV$ has a
subsequence converging in $L^1$, and the $L^1$ consistency condition identifies the
limit as $f$. The strict convergence upgrade follows from the lower semicontinuity of
the total variation functional and is the stronger result one would want for
applications requiring convergence of the total size of gradient information rather
than just its distribution.

The parallel between Lemmas~\ref{lem:sobolev_consistency}
and~\ref{lem:bv_consistency} is worth noting explicitly. In both cases the argument
has the same structure: uniform boundedness in a strong space gives compactness,
consistency in a weak space identifies the limit, and the combination gives
convergence of the derivative in the appropriate topology. The difference is in the
topology of the conclusion---strong $L^p$ convergence in the Sobolev case, weak-$*$
convergence of measures in the BV case---and this difference is fundamental rather
than technical. It reflects the fact that smooth estimators carry gradient
information in a stronger sense than piecewise constant estimators, which is itself a
reflection of the different inductive biases of the estimator classes.

\begin{lemma}[Hadamard Differentiability of the AME Functional]
	\label{lem:hadamard}
	Let $P$ be the marginal distribution of the covariates with density bounded away
	from zero and infinity on $\Omega$, and let $f \in \Wonep$ for $p > d$. The
	functional:
	\begin{equation}
		T: \mathcal{F} \to \R^d, \quad
		T(f) = \int_\Omega \nabla f(x) \, dP(x)
	\end{equation}
	is Hadamard differentiable at $f$ tangentially to $C(\Omega)^d$, with derivative:
	\begin{equation}
		T'_f(h) = \int_\Omega \nabla h(x) \, dP(x)
	\end{equation}
\end{lemma}

Lemma~\ref{lem:hadamard} is the keystone result of the paper. It establishes that
the AME functional varies smoothly when the underlying function is perturbed, in
precisely the sense required by the functional delta method
\citep[Chapter~3.9]{vandervaart1996weak}. Hadamard differentiability tangentially to
$C(\Omega)^d$ means that for any sequence $t_n \to 0$ and $h_n \to h$ in
$C(\Omega)^d$:
\begin{equation}
	\frac{T(f + t_n h_n) - T(f)}{t_n} \to T'_f(h) = \int_\Omega \nabla h \, dP
\end{equation}
This is the precise smoothness condition that allows the bootstrap to consistently
estimate the sampling distribution of $T(\fn)$. Intuitively, if small perturbations
to the estimator produce proportionally small changes in the AME in a controlled way,
then the bootstrap perturbations induced by resampling the data will correctly
characterize the variability of the AME estimate.

The condition $p > d$ ensures the Sobolev embedding $W^{1,p} \hookrightarrow
C^{0,\alpha}$ holds, which is required to control the continuity of $\nabla f$ and
make the integral well defined for perturbed functions. This is the one assumption
that is stronger than what the gradient consistency results require, and it is the
price of having a general bootstrap validity result.

The proof proceeds by verifying the Hadamard derivative formula directly and checking
the continuity conditions. Lemmas~\ref{lem:sobolev_existence}
through~\ref{lem:bv_consistency} establish that the gradient of the estimator
converges in the appropriate topology; Lemma~\ref{lem:hadamard} then shows that this
convergence is smooth enough to justify bootstrap inference. The two parts of the
theory --- gradient consistency and bootstrap validity --- are therefore not
independent results but facets of the same underlying smoothness property of the AME
functional.

\subsection{Main Theorems}

The main theorems apply the lemmas to each estimator class. In each case the argument
has two components: verifying that the estimator satisfies the well-behaved condition,
and citing the appropriate lemma pair. The first component is the estimator-specific
contribution; the second is the functional analytic machinery established above.

\begin{theorem}[Gradient Consistency for Neural Networks]
	\label{thm:neural_nets}
	Let $\fn$ be a neural network estimator with smooth activations $\sigma \in C^2$,
	bounded weight norms $\sup_n \prod_\ell \norm{W_\ell^{(n)}} < \infty$, and
	$\norm{\fn - f}_{L^p} \xrightarrowp 0$ for $p > d$. Then $\fn$ is well-behaved in
	the Sobolev sense and:
	\begin{equation}
		\norm{\nabla \fn - \nabla f}_{L^p} \xrightarrowp 0
	\end{equation}
\end{theorem}

The conditions of Theorem~\ref{thm:neural_nets} are natural and verifiable. Smooth
activations --- sigmoid, tanh, GELU --- are standard in practice and ensure the network
function is twice differentiable. The bounded weight norm condition formalizes the
effect of $L^2$ regularization during training, which prevents the network weights from
growing without bound. The $L^p$ consistency condition is the standard requirement that
the network is learning the true function. \citet{hornik1990universal} establish that
multilayer feedforward networks with smooth activation functions are capable of
approximating an arbitrary function and its derivatives to arbitrary accuracy, providing
the foundation for the consistency condition and establishing that the function class is
rich enough to contain the estimators we consider. The condition $p > d$ is required by
the pointwise differentiability theorem of \citet[Theorem~6.5]{evans2015measure} and
is satisfied for all $p$ larger than the covariate dimension --- a mild regularity
condition that is implicit in the classical Sobolev embedding theory.

\begin{proof}[Proof sketch]
	The chain rule applied through the network layers gives $\fn \in \Wtwop$ with:
	\begin{equation}
		\norm{\fn}_{W^{2,p}} \lesssim \prod_\ell \norm{W_\ell^{(n)}}^2 \cdot
		\norm{\sigma''}_\infty
	\end{equation}
	which is uniformly bounded by the weight norm assumption. The $L^p$ consistency
	condition and Lemma~\ref{lem:sobolev_consistency} then give the result. Full proof
	in the Supplementary Material \citep{parker2026supp}.
\end{proof}

\begin{remark}
	Theorem~\ref{thm:neural_nets} is stated for smooth activations. ReLU networks
	require a modified argument because ReLU is not $C^2$ and the chain rule calculation
	does not apply directly. Extension to ReLU networks is possible through a
	distributional argument but complicates the presentation and is omitted here.
\end{remark}

\begin{theorem}[Gradient Consistency for Gradient Boosted Models]
	\label{thm:boosting}
	Let $\fn = \hat{F}_m = (I - (I-S)^{m+1})Y$ be an $L_2$Boost estimator with
	symmetric linear learner $S$ satisfying the eigenvalue condition $0 < \lambda_k \leq
	1$ for all $k$, and with base learners $h_m$ that are either smooth ($h_m \in
	W^{2,p}$) or tree stumps. Suppose $\norm{\fn - f}_{L^1} \xrightarrowp 0$. Then:
	\begin{enumerate}
		\item[(i)] If base learners are smooth, $\fn$ is well-behaved in the Sobolev
		sense and $\norm{\nabla \fn - \nabla f}_{L^p} \xrightarrowp 0$.
		\item[(ii)] If base learners are tree stumps, $\fn$ is well-behaved in the
		BV sense and $D\fn \wstar Df$ in $\cM$.
	\end{enumerate}
\end{theorem}

Theorem~\ref{thm:boosting} covers both cases that arise in practice and illustrates
how the function space framework adapts to the base learner. XGBoost and
LightGBM --- the implementations used in Section~3 --- use shallow trees as base
learners and therefore fall under case~(ii). The two cases in the theorem are not
merely technical alternatives --- they reflect a genuine difference in what the
gradient of the estimator is. For smooth base learners the gradient is an ordinary
function and $L^p$ convergence is available. For tree base learners the gradient is a
measure and only weak-$*$ convergence can be established. The average marginal effects
are well defined and consistent in both cases, but the sense in which they are
consistent differs.

The eigenvalue condition $0 < \lambda_k \leq 1$ is the key regularity condition on
the base learner. \citet{buhlmann2003boosting} impose this condition explicitly in
their Theorem~1 and establish that under it the bias decays exponentially fast while
the variance increases only by exponentially diminishing amounts --- the exponential
bias--variance trade-off that makes $L_2$Boost analytically tractable. The condition
holds for smoothing spline learners and, generically, for tree stump learners whose
splits have positive but incomplete explanatory power, which is the typical case in
practice. Learners with any eigenvalue equal to zero are degenerate in the sense of
Corollary~1 of \citet{buhlmann2003boosting}: $\mathcal{B}_m \equiv S$ for all $m$ and
boosting does not improve on the base learner. The functional gradient descent
perspective of \citet{mason2000boosting} provides the conceptual framework for
understanding the boosting path as gradient descent in function space.

\begin{proof}[Proof sketch]
	For smooth base learners the $\Wtwop$ bound follows from linearity of the estimator:
	\begin{equation}
		\norm{\fn}_{W^{2,p}} \leq \sum_m \eta_m \norm{h_m}_{W^{2,p}} < \infty
	\end{equation}
	For tree stump base learners the BV bound uses the eigenvalue structure of
	$L_2$Boost. By Proposition~1 of \citet{buhlmann2003boosting}, the residual at step
	$m$ is $u_m = (I-S)^m Y$. By Proposition~2, under the eigenvalue condition $0 <
	\lambda_k \leq 1$, Theorem~1 of \citet{buhlmann2003boosting} gives exponential decay
	of residuals:
	\begin{equation}
		\norm{u_m}_{L^2} \leq C\rho^m, \quad \rho = \max_k(1-\lambda_k) < 1
	\end{equation}
	Each stump's BV norm is controlled by its residual magnitude via Lemma~\ref{lem:bv_existence},
	and the total BV norm of the ensemble satisfies:
	\begin{equation}
		\mathbb{E}[|D\fn|(\Omega)] \leq C\sum_{m=1}^\infty \rho^m = \frac{C\rho}{1-\rho} < \infty
	\end{equation}
	The geometric series converges because $\rho < 1$, giving uniform BV boundedness
	without any step size condition. Lemma~\ref{lem:bv_consistency} then gives
	$D\fn \xrightarrow{w*} Df$. Full proof in the Supplementary Material \citep{parker2026supp}.
\end{proof}

\begin{theorem}[Gradient Consistency for Random Forests]
	\label{thm:forests}
	Let $\fn$ be a Mondrian random forest estimator with budget $\lambda_n \asymp
	n^{1/(d+2)}$, $B \gg n^{2(d-\alpha)/(d+2)}$ trees where $\alpha = 1 - d/p$, $f$
	Lipschitz on $[0,1]^d$, and $f \in \Wonep$ for $p > d$. Then $\fn$ is well-behaved
	in the BV sense and:
	\begin{equation}
		D\fn \wstar Df \quad \text{in } \cM
	\end{equation}
\end{theorem}

Theorem~\ref{thm:forests} is the most technically demanding result in the paper.
Random forests are genuinely different from neural networks and boosted models in a
way that makes gradient consistency harder to establish --- the partition structure is
random and the response averages within leaves depend on the data in a complex way.
The restriction to Mondrian forests, where splits are made independently of the data,
is the key simplification that makes the proof tractable. For Mondrian forests the
partition geometry and the response variation can be controlled separately, whereas for
Breiman forests with data-dependent splits the two are coupled in a way that resists
clean analysis. The $L^2$ consistency result under the budget $\lambda_n \asymp
n^{1/(d+2)}$ follows from \citet{mourtada2017mondrian}, who establish the minimax rate
$O(n^{-2/(d+2)})$ for Lipschitz regression functions, and $L^1$ consistency follows
by Cauchy--Schwarz on the bounded domain $\Omega$.

The budget condition $\lambda_n \asymp n^{1/(d+2)}$ and the tree count condition $B
\gg n^{2(d-\alpha)/(d+2)}$ are natural and interpretable. The budget determines the
approximation rate --- balancing the bias from coarse partitions against the variance
from fine partitions at the minimax optimal rate. The tree count condition ensures that
the forest average has sufficient cancellation across trees to control the total
variation; it is mild in practice, satisfied by standard forest sizes for moderate
dimension and sample size.

\begin{proof}[Proof sketch]
	Each tree is piecewise constant and therefore in $BV$ by
	Lemma~\ref{lem:bv_existence}. The key difficulty is that individual tree BV norms
	grow with $n$, so the triangle inequality applied to the forest average is too crude.
	Instead we work with the second moment of the total variation. Writing $D\fn =
	\frac{1}{B}\sum_b DT_b$:
	\begin{equation}
		\mathbb{E}[|D\fn|(\Omega)] \leq \sqrt{\mathbb{E}[|D\fn|(\Omega)^2]}
		= \sqrt{\frac{1}{B^2}\mathbb{E}\!\left[\left|\sum_b DT_b\right|(\Omega)^2\right]}
	\end{equation}
	The cross terms between trees $b \neq b'$ vanish by independence of the Mondrian
	partitions and the zero-mean condition on the jumps. The diagonal terms are bounded
	using the cell diameter result $\mathbb{E}[D_{\lambda_n}(x)^2] \leq 4d/\lambda_n^2$
	of \citet{roy2009mondrian} and the split count bound $\mathbb{E}[K_{\lambda_n}]
	\leq (e(\lambda_n+1))^d$ of \citet{mourtada2017mondrian}, giving:
	\begin{equation}
		\mathbb{E}[|D\fn|(\Omega)^2]
		= \frac{1}{B}\,\mathbb{E}[|DT_1|(\Omega)^2]
		\leq \frac{1}{B}\,O(n^{2(d-\alpha)/(d+2)})
	\end{equation}
	Under the condition $B \gg n^{2(d-\alpha)/(d+2)}$ this is uniformly bounded.
	Lemma~\ref{lem:bv_consistency} then gives $D\fn \xrightarrow{w*} Df$ in $\cM$.
	Full proof in the Supplementary Material \citep{parker2026supp}.
\end{proof}

\begin{remark}
	Extension of Theorem~\ref{thm:forests} to Breiman forests with data-dependent splits
	requires additional control over the interaction between partition geometry and
	response variation. The data-dependent splitting rule couples these two quantities in
	a way that makes the cross-term cancellation argument unavailable --- splits are
	chosen precisely to maximize response differences across boundaries, which is the
	worst case for BV control. This extension is an important direction for future work
	and is discussed further in Section~\ref{sec:conclusion}.
\end{remark}

\subsection{Bootstrap Validity}

The gradient consistency results established above, combined with the Hadamard
differentiability of the AME functional established in Lemma~\ref{lem:hadamard},
yield bootstrap validity for average marginal effects across all three estimator
classes. The proof strategy in each case is the same: the gradient consistency
theorem establishes that $\fn$ converges to $f$ in the appropriate function space,
the rate condition verifies the speed of that convergence, and
Lemma~\ref{lem:hadamard} together with the functional delta method of
\citet[Theorem~3.9.4]{vandervaart1996weak} then gives bootstrap validity.

The three corollaries below make this argument explicit for each estimator class.
The rate conditions differ across classes in a way that reflects fundamental
differences in the estimators. For neural networks in the overparameterized regime
the parametric rate $n^{-1/2}$ is achievable and the rate condition is stated as
an explicit assumption that we conjecture is provable under the neural tangent
kernel approximation \citep{jacot2018ntk}. For gradient boosted models the
parametric rate holds under the eigenvalue condition of \citet{buhlmann2003boosting},
which governs the exponential decay of residuals along the boosting path and implies
convergence at the parametric rate. For Mondrian forests the rate is
dimension-dependent at $n^{-1/(d+2)}$, reflecting the minimax lower bound for
nonparametric estimation of Lipschitz functions --- a fundamental limitation of the
problem rather than of the method.

\begin{corollary}[Bootstrap Validity for Neural Network AMEs]
	\label{cor:bootstrap_nn}
	Under the conditions of Theorem~\ref{thm:neural_nets}, suppose additionally that
	$\sigma \in C^3$ and that:
	\begin{equation}
		\sqrt{n}(\fn - f) = O_p(1) \quad \text{in } \Wonep
		\tag{RC}
	\end{equation}
	Then the bootstrap distribution of:
	\begin{equation}
		\sqrt{n}\left(\frac{1}{n}\sum_{i=1}^n \nabla \hat{f}_n^*(x_i) -
		\frac{1}{n}\sum_{i=1}^n \nabla \fn(x_i)\right)
	\end{equation}
	consistently estimates the sampling distribution of
	$\sqrt{n}\left(\frac{1}{n}\sum_{i=1}^n \nabla \fn(x_i) - T(f)\right)$.
\end{corollary}

The condition $\sigma \in C^3$ rather than $C^2$ is required because the Hadamard
differentiability argument in Lemma~\ref{lem:hadamard} needs one additional derivative
beyond what Theorem~\ref{thm:neural_nets} requires --- specifically to control the
smoothness of the bootstrap perturbations of the network weights as they propagate
through the gradient of the network function.

The rate condition (RC) is stated as an explicit assumption rather than a proved
result. We conjecture it holds in the overparameterized regime where the neural
tangent kernel approximation applies. In this regime the network training dynamics
are approximately linear in the weights \citep{jacot2018ntk}, the weight estimates
converge at the parametric rate $n^{-1/2}$, and propagating this through the network
function via the $C^3$ chain rule calculation gives $O_p(n^{-1/2})$ convergence in
$\Wonep$. A complete proof requires quantitative finite-width NTK approximation
results and is left as future work. Unlike the Mondrian forest case where the
dimension-dependent rate is a fundamental consequence of the nonparametric minimax
lower bound, we expect the parametric rate to be achievable for neural networks
without a dimension penalty, making (RC) a provable rather than fundamental
restriction.

\begin{corollary}[Bootstrap Validity for Gradient Boosted AMEs]
	\label{cor:bootstrap_boost}
	Under the conditions of Theorem~\ref{thm:boosting} with eigenvalue condition
	$0 < \lambda_k \leq 1$, suppose additionally that $\sqrt{n}(\fn - f) = O_p(1)$
	in the appropriate norm. Then the bootstrap distribution of:
	\begin{equation}
		\sqrt{n}\left(\frac{1}{n}\sum_{i=1}^n D\hat{f}_n^*(x_i) -
		\frac{1}{n}\sum_{i=1}^n D\fn(x_i)\right)
	\end{equation}
	consistently estimates the sampling distribution of
	$\sqrt{n}\left(\frac{1}{n}\sum_{i=1}^n D\fn(x_i) - T(f)\right)$.
\end{corollary}

The rate condition for gradient boosting follows from the exponential residual decay
established in the proof of Theorem~\ref{thm:boosting}. By Proposition~1 and
Theorem~1 of \citet{buhlmann2003boosting}, under the eigenvalue condition $0 <
\lambda_k \leq 1$, the residuals decay geometrically as $\|u_m\|_{L^2} \leq C\rho^m$
for $\rho = \max_k(1-\lambda_k) < 1$. This geometric decay implies convergence of
the boosting estimator at a rate sufficient for the parametric $n^{-1/2}$ rate,
since the bias vanishes exponentially fast and the variance is controlled by the
eigenvalue structure of $S$. The bootstrap validity argument then adapts the
framework of \citet{horowitz2000bootstrap} for smooth functionals of empirical
processes to the sequential structure of the boosting path, using the geometric
decay of residuals to control the dependence across iterations.

\begin{corollary}[Bootstrap Validity for Mondrian Forest AMEs]
	\label{cor:bootstrap_rf}
	Under the conditions of Theorem~\ref{thm:forests}, the bootstrap distribution of:
	\begin{equation}
		n^{1/(d+2)}\left(\frac{1}{n}\sum_{i=1}^n D\hat{f}_n^*(x_i) -
		\frac{1}{n}\sum_{i=1}^n D\fn(x_i)\right)
	\end{equation}
	consistently estimates the sampling distribution of:
	\begin{equation}
		n^{1/(d+2)}\left(\frac{1}{n}\sum_{i=1}^n D\fn(x_i) - T(f)\right)
	\end{equation}
	In particular, when $d \leq 2$ the scaling $n^{1/(d+2)}$ achieves the parametric
	rate $n^{1/2}$ and the bootstrap confidence intervals are asymptotically exact at
	the parametric rate.
\end{corollary}

The Mondrian forest case differs from the neural network and gradient boosting cases
in that bootstrap validity holds at a dimension-dependent rate rather than the
parametric rate $n^{1/2}$ universally. The rate $n^{1/(d+2)}$ is the minimax optimal
rate for Lipschitz regression functions on $[0,1]^d$, established by
\citet{mourtada2017mondrian} with the budget sequence $\lambda_n \asymp n^{1/(d+2)}$.
This rate achieves $n^{1/2}$ when $d \leq 2$, and is slower than $n^{1/2}$ for $d
\geq 3$ --- a consequence of the curse of dimensionality that is fundamental to
nonparametric estimation rather than a limitation of the present framework. The
bootstrap remains valid at the actual convergence rate in all dimensions, consistently
estimating the sampling distribution of the AME estimator regardless of dimension.

The budget condition $\lambda_n \asymp n^{1/(d+2)}$ plays two roles. First it ensures
the minimax rate is achieved, balancing bias from coarse partitions against variance
from fine partitions. Second it determines the rate at which the BV norm of the
forest average is controlled via the second moment argument in the proof of
Theorem~\ref{thm:forests} --- specifically the cell diameter bound
$\mathbb{E}[D_{\lambda_n}(x)^2] \leq 4d/\lambda_n^2 = O(n^{-2/(d+2)})$ of
\citet{roy2009mondrian} and the split count bound $\mathbb{E}[K_{\lambda_n}] \leq
(e(\lambda_n+1))^d = O(n^{d/(d+2)})$ of \citet{mourtada2017mondrian}.

The independence of the Mondrian partition from the data --- the same property that
made the gradient consistency proof tractable in Theorem~\ref{thm:forests} --- also
simplifies the bootstrap validity argument. Under bootstrap resampling the partition
$\mathcal{P}_b$ is unchanged across replicates and only the response values within
fixed leaves are perturbed. The conditional independence of subtrees established by
\citet{roy2009mondrian} then controls the dependence structure of the bootstrap
gradient measures across leaves, completing the bootstrap validity argument via
Lemma~\ref{lem:hadamard} and \citet[Theorem~3.9.4]{vandervaart1996weak}. The
connection to the infinitesimal jackknife results of \citet{wager2014confidence}
provides additional scaffolding: their variance estimation framework for forest
estimators is essentially a first-order bootstrap approximation, and the Hadamard
differentiability established in Lemma~\ref{lem:hadamard} shows that this
approximation is asymptotically exact at the Mondrian convergence rate.

\begin{remark}
	The three corollaries establish bootstrap validity for neural networks in the
	overparameterized regime, gradient boosted models, and Mondrian forests. The rate
	conditions reflect the fundamental structure of each estimator class: the parametric
	rate for neural networks and gradient boosting, and the dimension-dependent minimax
	rate for Mondrian forests. Bootstrap validity for Breiman forests with
	data-dependent splits remains open for the same reason that gradient consistency
	for Breiman forests is harder --- the data-dependent partition geometry couples the
	bootstrap perturbations to the partition structure in a way that prevents the cross-
	term cancellation argument and resists clean analysis. This is an important
	direction for future work.
\end{remark}

\section{Empirical Examples}

Across the three empirical applications we fit four model classes: logistic or linear regression as the parametric baseline, a random forest, XGBoost, and a feedforward neural network. The random forest uses 200 trees with a maximum depth of 8. XGBoost uses 200 boosting rounds with a maximum depth of 4 and a learning rate of 0.05. The neural network architecture and training schedule depend on the outcome type. For binary classification, the network is a feedforward architecture with two hidden layers of 64 and 32 units trained with the Adam optimizer at a learning rate of $10^{-3}$, with early stopping on a 10\% validation split with patience of 20 epochs. For continuous regression, the network is deeper --- four hidden layers of 256, 128, 64, and 32 units --- trained at a learning rate of $3 \times 10^{-4}$ with patience of 50 epochs, reflecting the larger capacity required to fit the continuous outcome surface and the longer training horizon needed for the optimization to converge. For regression outcomes, the network includes the input normalization layer and target standardization wrapper described in Section~2 to ensure scale-invariant gradient updates and predictions on the natural outcome scale. All models are trained on the full sample for each application, with the bootstrap procedure described in Section~2 producing standard errors via 200 replicates. Where applicable, predicted probabilities for classification models are extracted from \texttt{predict\_proba(X)[:,1]} as discussed in Section~2.

\subsection{UCI Adult Income}

The first dataset is the UCI Adult Income dataset \citep{UCI, kohavi1996}, a widely used benchmark for binary classification in the social sciences. We combine the training and test sets for a total of 48,842 observations. The outcome is whether an individual earns more than \$50,000 per year. Features are organized into three panels: demographics (age, sex, years of education), work characteristics (hours worked per week, occupation group, work class), and financial variables (capital gains, capital losses, marital status). The specification search considers logistic regression and XGBoost across four specifications, allowing the method to be evaluated both where the parametric model is approximately correctly specified and where the black-box learner offers efficiency gains.

\begin{table*}[htbp]
	\centering
	\small
	
	\resizebox{\textwidth}{!}{%
	\begin{threeparttable}
		\caption{UCI Adult Income: Specification Search. Average marginal effects on P(Income $>$ \$50k).}
		\label{tab:adult_spec_search}
		
		\begin{tabular}{lrrrrrrrr}
			\toprule
			& \multicolumn{4}{c}{Logistic} & \multicolumn{4}{c}{XGBoost} \\
			\cmidrule(lr){2-5} \cmidrule(lr){6-9} 
			& Spec A & A+B & A+C & A+B+C & Spec A & A+B & A+C & A+B+C \\
			\midrule
			\multicolumn{9}{l}{\textit{Panel A: Demographics}} \\
			\quad age & 0.0060$^{***}$ & 0.0060$^{***}$ & 0.0024$^{***}$ & 0.0027$^{***}$ & 0.0059$^{***}$ & 0.0059$^{***}$ & 0.0029$^{***}$ & 0.0031$^{***}$ \\
			& (0.0001) & (0.0001) & (0.0001) & (0.0001) & (0.0002) & (0.0001) & (0.0001) & (0.0001) \\
			\quad female & -0.1730$^{***}$ & -0.1648$^{***}$ & -0.0275$^{***}$ & -0.0269$^{***}$ & -0.1592$^{***}$ & -0.1446$^{***}$ & -0.0138$^{***}$ & -0.0122$^{***}$ \\
			& (0.0034) & (0.0034) & (0.0041) & (0.0041) & (0.0032) & (0.0032) & (0.0028) & (0.0032) \\
			\quad education\_num & 0.0523$^{***}$ & 0.0393$^{***}$ & 0.0400$^{***}$ & 0.0306$^{***}$ & 0.0610$^{***}$ & 0.0452$^{***}$ & 0.0405$^{***}$ & 0.0312$^{***}$ \\
			& (0.0006) & (0.0007) & (0.0006) & (0.0006) & (0.0023) & (0.0017) & (0.0011) & (0.0011) \\
			\midrule
			\multicolumn{9}{l}{\textit{Panel B: Work Characteristics}} \\
			\quad hours\_per\_week & -- & 0.0048$^{***}$ & -- & 0.0033$^{***}$ & -- & 0.0084$^{***}$ & -- & 0.0032 \\
			&  & (0.0001) &  & (0.0001) &  & (0.0028) &  & (0.0025) \\
			\quad government & -- & 0.0029 & -- & 0.0033 & -- & -0.0067$^{*}$ & -- & -0.0003 \\
			&  & (0.0051) &  & (0.0045) &  & (0.0037) &  & (0.0026) \\
			\quad self\_employed & -- & -0.0077 & -- & -0.0268$^{***}$ & -- & 0.0113$^{**}$ & -- & -0.0055$^{*}$ \\
			&  & (0.0050) &  & (0.0044) &  & (0.0054) &  & (0.0033) \\
			\quad white\_collar & -- & 0.0666$^{***}$ & -- & 0.0541$^{***}$ & -- & 0.0466$^{***}$ & -- & 0.0474$^{***}$ \\
			&  & (0.0056) &  & (0.0051) &  & (0.0026) &  & (0.0017) \\
			\quad blue\_collar & -- & -0.0264$^{***}$ & -- & -0.0270$^{***}$ & -- & -0.0446$^{***}$ & -- & -0.0251$^{***}$ \\
			&  & (0.0052) &  & (0.0048) &  & (0.0041) &  & (0.0034) \\
			\quad service & -- & -0.0697$^{***}$ & -- & -0.0427$^{***}$ & -- & -0.0807$^{***}$ & -- & -0.0343$^{***}$ \\
			&  & (0.0069) &  & (0.0066) &  & (0.0049) &  & (0.0058) \\
			\midrule
			\multicolumn{9}{l}{\textit{Panel C: Financial}} \\
			\quad capital\_gain & -- & -- & 0.0000$^{***}$ & 0.0000$^{***}$ & -- & -- & -0.0001$^{***}$ & -0.0001$^{***}$ \\
			&  &  & (0.0000) & (0.0000) &  &  & (0.0000) & (0.0000) \\
			\quad capital\_loss & -- & -- & 0.0001$^{***}$ & 0.0001$^{***}$ & -- & -- & -0.0001$^{***}$ & -0.0000$^{***}$ \\
			&  &  & (0.0000) & (0.0000) &  &  & (0.0000) & (0.0000) \\
			\quad married & -- & -- & 0.2854$^{***}$ & 0.2752$^{***}$ & -- & -- & 0.2383$^{***}$ & 0.2357$^{***}$ \\
			&  &  & (0.0042) & (0.0041) &  &  & (0.0035) & (0.0036) \\
			\midrule
			McFadden $R^2$ & 0.199 & 0.233 & 0.376 & 0.400 & 0.257 & 0.296 & 0.467 & 0.492 \\
			Accuracy & 0.797 & 0.806 & 0.841 & 0.847 & 0.805 & 0.818 & 0.864 & 0.872 \\
			\bottomrule
		\end{tabular}
		\begin{tablenotes}
			\footnotesize
			\item \textit{Notes:} $^{***}$p$<$0.01, $^{**}$p$<$0.05, $^{*}$p$<$0.10. Standard errors from nonparametric bootstrap (200 replicates) in parentheses.
		\end{tablenotes}
	\end{threeparttable}
	}%
\end{table*}

Table~\ref{tab:adult_spec_search} presents the specification search results for the UCI Adult Income dataset. The demographic panel tells a consistent story across all specifications and both model classes. Age and education have positive effects on the probability of earning above \$50,000, and the female indicator is large and negative, falling from around -\$0.16 to $-$0.03 as financial and work controls are added --- a pattern of omitted variable bias that is clearly legible in the AME framework and consistent across logistic regression and XGBoost. The standard errors are tight throughout, reflecting the large sample size, and the XGBoost standard errors are comparable to or smaller than those from logistic regression despite the additional model complexity. This is the efficiency gain the framework is designed to produce: when the black-box model fits better, as reflected in the higher McFadden $R^2$ across all specifications, the bootstrap standard errors narrow rather than widen. The financial panel illustrates a limitation of the capital gains and losses variables --- the AMEs are near zero in magnitude because these variables are highly skewed, and the units are raw dollars, so the displayed coefficients reflect a one-dollar change. The sign reversal on capital gains between logistic regression and XGBoost in specification A+C warrants attention and likely reflects nonlinearity that the logistic model cannot capture.

\begin{table*}[htbp]
	\centering
	\small
	\begin{threeparttable}
		\caption{UCI Adult Income: Model Comparison. Average marginal effects on P(Income $>$ \$50k), full specification (A+B+C).}
		\label{tab:adult_model_comparison}
		\begin{tabular}{lrrrr}
			\toprule
			Feature & Logistic & Random Forest & XGBoost & Neural Net \\
			\midrule
			age & 0.0027$^{***}$ & 0.0042$^{***}$ & 0.0031$^{***}$ & 0.0035$^{***}$ \\
			& (0.0001) & (0.0001) & (0.0001) & (0.0006) \\
			female & -0.0269$^{***}$ & -0.0218$^{***}$ & -0.0122$^{***}$ & -0.0085 \\
			& (0.0041) & (0.0001) & (0.0032) & (0.0084) \\
			education\_num & 0.0306$^{***}$ & 0.0294$^{***}$ & 0.0312$^{***}$ & 0.0238$^{***}$ \\
			& (0.0006) & (0.0002) & (0.0011) & (0.0025) \\
			\midrule
			hours\_per\_week & 0.0033$^{***}$ & 0.0070$^{***}$ & 0.0032 & 0.0028$^{***}$ \\
			& (0.0001) & (0.0001) & (0.0025) & (0.0006) \\
			government & 0.0033 & 0.0051$^{***}$ & -0.0003 & -0.0065 \\
			& (0.0045) & (0.0001) & (0.0026) & (0.0090) \\
			self\_employed & -0.0268$^{***}$ & 0.0032$^{***}$ & -0.0055$^{*}$ & -0.0068 \\
			& (0.0044) & (0.0001) & (0.0033) & (0.0066) \\
			white\_collar & 0.0541$^{***}$ & 0.0600$^{***}$ & 0.0474$^{***}$ & 0.0243$^{***}$ \\
			& (0.0051) & (0.0002) & (0.0017) & (0.0065) \\
			blue\_collar & -0.0270$^{***}$ & -0.0240$^{***}$ & -0.0251$^{***}$ & -0.0502$^{***}$ \\
			& (0.0048) & (0.0001) & (0.0034) & (0.0081) \\
			service & -0.0427$^{***}$ & -0.0176$^{***}$ & -0.0343$^{***}$ & -0.0452$^{***}$ \\
			& (0.0066) & (0.0001) & (0.0058) & (0.0088) \\
			\midrule
			capital\_gain & 0.0000$^{***}$ & 0.0000$^{***}$ & -0.0001$^{***}$ & 0.0001$^{***}$ \\
			& (0.0000) & (0.0000) & (0.0000) & (0.0000) \\
			capital\_loss & 0.0001$^{***}$ & -0.0000$^{***}$ & -0.0000$^{***}$ & 0.0037$^{***}$ \\
			& (0.0000) & (0.0000) & (0.0000) & (0.0006) \\
			married & 0.2752$^{***}$ & 0.2446$^{***}$ & 0.2357$^{***}$ & 0.2046$^{***}$ \\
			& (0.0041) & (0.0006) & (0.0036) & (0.0196) \\
			\midrule
			McFadden $R^2$ & 0.400 & 0.450 & 0.492 & 0.418 \\
			Accuracy & 0.847 & 0.863 & 0.872 & 0.853 \\
			\bottomrule
		\end{tabular}
		\begin{tablenotes}
			\footnotesize
			\item \textit{Notes:} $^{***}$p$<$0.01, $^{**}$p$<$0.05, $^{*}$p$<$0.10. Standard errors from nonparametric bootstrap (200 replicates) in parentheses.
		\end{tablenotes}
	\end{threeparttable}
\end{table*}

Table~\ref{tab:adult_model_comparison} compares all four model classes on the full specification. The most important finding is the stability of the core demographic results. The education AME is 0.031 for logistic regression, 0.029 for random forest, 0.031 for XGBoost, and 0.024 for the neural network --- a range of less than one percentage point across models that differ enormously in their functional form. The married AME is similarly stable, ranging from 0.205 to 0.275. This stability is not guaranteed by the method --- it reflects genuine agreement across models about the underlying conditional expectation function --- and it is the kind of cross-model validation that the AME framework makes possible precisely because it produces output on a common scale. The female indicator shows more variation, from $-$0.027 for logistic regression to $-$0.009 for the neural network, and is not significant for the neural network, suggesting some sensitivity to model specification for this variable. The random forest standard errors are notably smaller than those from other models for most variables, which is a consequence of the ensemble averaging in that estimator producing a smoother prediction surface and therefore more stable finite differences across bootstrap replicates.

\begin{sidewaystable}[htbp]
	\centering
	\small
	\begin{threeparttable}
		\caption{UCI Adult Income: Method Comparison. AME, SHAP, and PDP estimates for P(Income $>$ \$50k), full specification (A+B+C).}
		\label{tab:adult_method_comparison}
		\begin{tabular}{lrrrrrrrrrrrr}
			\toprule
			& \multicolumn{3}{c}{Logistic} & \multicolumn{3}{c}{Random Forest} & \multicolumn{3}{c}{XGBoost} & \multicolumn{3}{c}{Neural Net} \\
			\cmidrule(lr){2-4} \cmidrule(lr){5-7} \cmidrule(lr){8-10} \cmidrule(lr){11-13} 
			& AME & SHAP & PDP & AME & SHAP & PDP & AME & SHAP & PDP & AME & SHAP & PDP \\
			\midrule
			age & 0.0027 & -0.0270 & 0.0028 & 0.0042 & -0.0068 & 0.0024 & 0.0031 & -0.0097 & 0.0036 & 0.0035 & 0.0044 & 0.0033 \\
			female & -0.0269 & 0.0022 & -0.0270 & -0.0218 & 0.0011 & -0.0218 & -0.0122 & 0.0015 & -0.0122 & -0.0085 & -0.0001 & -0.0222 \\
			education\_num & 0.0306 & -0.1186 & 0.0269 & 0.0294 & -0.0077 & 0.0137 & 0.0312 & -0.0093 & 0.0170 & 0.0238 & -0.0051 & 0.0207 \\
			\midrule
			hours\_per\_week & 0.0033 & -0.0691 & 0.0033 & 0.0070 & -0.0054 & 0.0020 & 0.0032 & -0.0080 & 0.0028 & 0.0028 & -0.0032 & 0.0028 \\
			government & 0.0033 & -0.0009 & 0.0026 & 0.0051 & -0.0003 & 0.0051 & -0.0003 & -0.0008 & -0.0003 & -0.0065 & 0.0002 & 0.0082 \\
			self\_employed & -0.0268 & 0.0016 & -0.0265 & 0.0032 & -0.0003 & 0.0032 & -0.0055 & -0.0002 & -0.0055 & -0.0068 & -0.0008 & -0.0045 \\
			white\_collar & 0.0541 & 0.0127 & 0.0538 & 0.0600 & 0.0031 & 0.0600 & 0.0474 & 0.0025 & 0.0474 & 0.0243 & 0.0007 & 0.0312 \\
			blue\_collar & -0.0270 & 0.0136 & -0.0273 & -0.0240 & 0.0010 & -0.0240 & -0.0251 & 0.0004 & -0.0251 & -0.0502 & -0.0008 & -0.0531 \\
			service & -0.0427 & -0.0027 & -0.0427 & -0.0176 & 0.0003 & -0.0176 & -0.0343 & 0.0005 & -0.0343 & -0.0452 & 0.0011 & -0.0558 \\
			\midrule
			capital\_gain & 0.0000 & -0.2627 & 0.0000 & 0.0000 & -0.0224 & 0.0000 & -0.0001 & -0.0264 & -0.0001 & 0.0001 & 0.0042 & 0.0000 \\
			capital\_loss & 0.0001 & 0.0027 & -- & -0.0000 & -0.0034 & -- & -0.0000 & -0.0062 & -- & 0.0037 & -0.0115 & -- \\
			married & 0.2752 & -0.0768 & 0.2749 & 0.2446 & -0.0064 & 0.2446 & 0.2357 & -0.0049 & 0.2357 & 0.2046 & 0.0149 & 0.1958 \\
			\bottomrule
		\end{tabular}
		\begin{tablenotes}
			\footnotesize
			\item \textit{Notes:} AME = Average Marginal Effect (marginfx). SHAP = mean signed SHAP value (GradientExplainer). PDP = slope of partial dependence plot. Full specification (A+B+C) only. No standard errors reported for SHAP and PDP.
		\end{tablenotes}
	\end{threeparttable}
\end{sidewaystable}

Table~\ref{tab:adult_method_comparison} compares AMEs against SHAP values and partial dependence plot slopes across all four models. The contrast is stark. The AME and PDP estimates are in close agreement throughout --- the married AME of 0.275 for logistic regression matches the PDP slope of 0.275, and similar agreement holds for the occupation indicators and hours per week. This agreement is reassuring: both methods are estimating derivatives of the prediction function, and their numerical proximity confirms that the finite difference approximation is well-behaved. The SHAP values, by contrast, are on a different scale and in several cases have the wrong sign relative to AME and PDP. The SHAP value for age is negative across all four models while the AME and PDP are both positive --- a direct contradiction that reflects the fact that mean signed SHAP values are not marginal effects and should not be interpreted as such. Crucially, no standard errors are reported for SHAP or PDP, so there is no basis for inference from those columns. A researcher using SHAP alone would have a point estimate of unknown precision with no ability to test whether age, education, or any other variable has a statistically distinguishable effect on the outcome.

\subsection{UCI Credit Default}

The second dataset is the UCI Default of Credit Card Clients dataset \citep{yeh2009, UCI}, comprising 30,000 observations of Taiwanese credit card holders. The outcome is whether a cardholder defaulted on their payment in the following month. Features are organized into demographics (age, sex, education, marital status), payment history (average payment status, months delayed), and bill and payment amounts (average bill amount, average payment amount, payment ratio). The specification search considers logistic regression and a TensorFlow neural network, providing a direct comparison between a classical parametric model and a deep learning approach on the same inferential task.

\begin{table*}[htbp]
	\centering
	\small
	\resizebox{\textwidth}{!}{%
		\begin{threeparttable}
			\caption{UCI Credit Card Default: Specification Search. Average marginal effects on P(Default).}
			\label{tab:credit_default_spec_search}
			\begin{tabular}{lrrrrrrrr}
				\toprule
				Variable & \multicolumn{4}{c}{Logistic} & \multicolumn{4}{c}{Neural Net} \\
				\cmidrule(lr){2-5} \cmidrule(lr){6-9} 
				& Spec A & A+B & A+C & A+B+C & Spec A & A+B & A+C & A+B+C \\
				\midrule
				\multicolumn{9}{l}{\textit{Panel A: Demographics}} \\
				\quad age & -0.0004 & 0.0002 & 0.0002 & 0.0002 & -0.0022$^{**}$ & 0.0001 & -0.0001 & 0.0007 \\
				& (0.0003) & (0.0003) & (0.0003) & (0.0002) & (0.0009) & (0.0008) & (0.0013) & (0.0005) \\
				\quad female & -0.0358$^{***}$ & -0.0207$^{***}$ & -0.0294$^{***}$ & -0.0198$^{***}$ & -0.0343$^{***}$ & -0.0204$^{***}$ & -0.0250$^{***}$ & -0.0135$^{***}$ \\
				& (0.0052) & (0.0045) & (0.0048) & (0.0045) & (0.0082) & (0.0064) & (0.0028) & (0.0028) \\
				\quad education & 0.0173$^{***}$ & 0.0058$^{**}$ & 0.0017 & 0.0002 & 0.0194$^{**}$ & 0.0166$^{***}$ & 0.0034 & 0.0078$^{***}$ \\
				& (0.0030) & (0.0028) & (0.0030) & (0.0028) & (0.0076) & (0.0046) & (0.0037) & (0.0013) \\
				\quad married & 0.0253$^{***}$ & 0.0186$^{***}$ & 0.0280$^{***}$ & 0.0220$^{***}$ & 0.0324$^{***}$ & 0.0145$^{**}$ & 0.0274$^{***}$ & 0.0192$^{***}$ \\
				& (0.0052) & (0.0049) & (0.0050) & (0.0050) & (0.0060) & (0.0058) & (0.0076) & (0.0049) \\
				\midrule
				\multicolumn{9}{l}{\textit{Panel B: Payment History}} \\
				\quad avg\_pay\_status & -- & -0.0161$^{***}$ & -- & -0.0256$^{***}$ & -- & 0.0231$^{***}$ & -- & 0.0270 \\
				&  & (0.0036) &  & (0.0042) &  & (0.0082) &  & (0.0299) \\
				\quad months\_delayed & -- & 0.0835$^{***}$ & -- & 0.0826$^{***}$ & -- & 0.1151$^{***}$ & -- & 0.1118$^{***}$ \\
				&  & (0.0019) &  & (0.0021) &  & (0.0165) &  & (0.0321) \\
				\midrule
				\multicolumn{9}{l}{\textit{Panel C: Bill and Payment Amounts}} \\
				\quad avg\_bill\_amt & -- & -- & -0.0000 & 0.0000 & -- & -- & 0.0000 & 0.0000 \\
				&  &  & (0.0000) & (0.0000) &  &  & (0.0000) & (0.0000) \\
				\quad avg\_pay\_amt & -- & -- & -0.0000$^{***}$ & -0.0000$^{***}$ & -- & -- & -0.0000 & -0.0000 \\
				&  &  & (0.0000) & (0.0000) &  &  & (0.0001) & (0.0001) \\
				\quad pay\_ratio & -- & -- & -0.1196$^{***}$ & -0.0454$^{***}$ & -- & -- & -0.1215$^{***}$ & 0.0009 \\
				&  &  & (0.0073) & (0.0072) &  &  & (0.0030) & (0.0029) \\
				\midrule
				McFadden $R^2$ & 0.003 & 0.130 & 0.031 & 0.137 & 0.011 & 0.152 & 0.060 & 0.171 \\
				Accuracy & 0.779 & 0.803 & 0.779 & 0.803 & 0.779 & 0.807 & 0.780 & 0.808 \\
				\bottomrule
			\end{tabular}
			\begin{tablenotes}
				\footnotesize
				\item \textit{Notes:} $^{***}$p$<$0.01, $^{**}$p$<$0.05, $^{*}$p$<$0.10. Standard errors from nonparametric bootstrap (200 replicates) in parentheses.
			\end{tablenotes}
		\end{threeparttable}
	}%
\end{table*}

Table~\ref{tab:credit_default_spec_search} presents the specification search. The dominant finding is the importance of payment history. Months delayed enters with an AME of 0.084 for logistic regression and 0.115 for the neural network in specification A+B, and remains large and significant in the full specification --- the largest effect in the table by a substantial margin. The demographic panel tells a different story from the Adult Income dataset: age is statistically indistinguishable from zero in nearly every specification for logistic regression, and the female indicator, while consistently negative and significant for both models, shrinks substantially once payment history is controlled for, falling from $-$\$0.036 to $-$0.020 for logistic regression and from $-$0.034 to $-$0.014 for the neural network. This pattern of attenuation is cleanly legible in the AME framework because the coefficients are on a common probability scale across specifications. The neural network and logistic regression produce broadly comparable AMEs throughout the demographic panel, with model fit improving steadily as features are added --- McFadden $R^2$ rising from 0.011 in specification A to 0.171 in the full specification for the neural network, slightly higher than logistic regression's 0.137. The bootstrap standard errors for the neural network are wider than those for logistic regression in the full specification, particularly for the payment history variables, reflecting the additional uncertainty introduced by hyperparameter selection on each replicate.

\begin{table*}[htbp]
	\centering
	\small
	\begin{threeparttable}
		\caption{UCI Credit Card Default: Model Comparison. Average marginal effects on P(Default), full specification (A+B+C).}
		\label{tab:credit_default_model_comparison}
		\begin{tabular}{lrrrr}
			\toprule
			Variable & Logistic & Random Forest & XGBoost & Neural Net \\
			\midrule
			age & 0.0002 & -0.0000 & -0.0004 & 0.0007 \\
			& (0.0002) & (0.0000) & (0.0006) & (0.0005) \\
			female & -0.0198$^{***}$ & -0.0062$^{***}$ & -0.0146$^{***}$ & -0.0135$^{***}$ \\
			& (0.0045) & (0.0000) & (0.0031) & (0.0028) \\
			education & 0.0002 & 0.0029$^{***}$ & 0.0030$^{*}$ & 0.0078$^{***}$ \\
			& (0.0028) & (0.0001) & (0.0016) & (0.0013) \\
			married & 0.0220$^{***}$ & 0.0054$^{***}$ & 0.0142$^{***}$ & 0.0192$^{***}$ \\
			& (0.0050) & (0.0000) & (0.0035) & (0.0049) \\
			\midrule
			avg\_pay\_status & -0.0256$^{***}$ & -0.0000 & 0.1027$^{***}$ & 0.0270 \\
			& (0.0042) & (0.0000) & (0.0225) & (0.0299) \\
			months\_delayed & 0.0826$^{***}$ & 0.0179$^{***}$ & 0.0209$^{***}$ & 0.1118$^{***}$ \\
			& (0.0021) & (0.0002) & (0.0009) & (0.0321) \\
			\midrule
			avg\_bill\_amt & 0.0000 & -0.0000$^{***}$ & -0.0000$^{*}$ & 0.0000 \\
			& (0.0000) & (0.0000) & (0.0000) & (0.0000) \\
			avg\_pay\_amt & -0.0000$^{***}$ & -0.0000$^{***}$ & -0.0000$^{***}$ & -0.0000 \\
			& (0.0000) & (0.0000) & (0.0000) & (0.0001) \\
			pay\_ratio & -0.0454$^{***}$ & -0.4306$^{***}$ & -0.4396$^{***}$ & 0.0009 \\
			& (0.0072) & (0.0058) & (0.1289) & (0.0029) \\
			\midrule
			McFadden $R^2$ & 0.137 & 0.210 & 0.194 & 0.171 \\
			Accuracy & 0.803 & 0.821 & 0.815 & 0.808 \\
			\bottomrule
		\end{tabular}
		\begin{tablenotes}
			\footnotesize
			\item \textit{Notes:} $^{***}$p$<$0.01, $^{**}$p$<$0.05, $^{*}$p$<$0.10. Standard errors from nonparametric bootstrap (200 replicates) in parentheses.
		\end{tablenotes}
	\end{threeparttable}
\end{table*}

Table~\ref{tab:credit_default_model_comparison} compares all four models on the full specification. The payment history variables again dominate, but the models disagree substantially on their magnitudes in a way that is informative rather than troubling. Months delayed is 0.083 for logistic regression, 0.018 for random forest, 0.021 for XGBoost, and 0.112 for the neural network --- a wide range that reflects genuine differences in how each model weighs this variable against the payment ratio, which random forest and XGBoost estimate at $-$0.431 and $-$0.440 respectively while logistic regression estimates only $-$0.045 and the neural network estimates near zero at 0.001. This is a case where the black-box tree-based models detect a strong nonlinear relationship between payment ratio and default probability, while the neural network instead loads that information onto months delayed, and logistic regression splits the effect between the two variables. The AME framework makes these disagreements visible and quantifiable in a way that point estimates from a single model would not. The random forest standard errors are again notably tight --- the female AME of $-$0.006 has a standard error of 0.000 to four decimal places --- reflecting the stability of the ensemble prediction surface. The neural network standard errors for the payment history variables are the widest in the table at around 0.030 to 0.032, which is the appropriate response of the bootstrap to a flexible model that distributes effects differently across replicates.

	\begin{sidewaystable}[htbp]
	\centering
	\small
	\begin{threeparttable}
		\caption{UCI Credit Card Default: Method Comparison. AME, SHAP, and PDP estimates for P(Default), full specification (A+B+C).}
		\label{tab:credit_default_method_comparison}
		\begin{tabular}{lrrrrrrrrrrrr}
			\toprule
			Variable & \multicolumn{3}{c}{Logistic} & \multicolumn{3}{c}{Random Forest} & \multicolumn{3}{c}{XGBoost} & \multicolumn{3}{c}{Neural Net} \\
			\cmidrule(lr){2-4} \cmidrule(lr){5-7} \cmidrule(lr){8-10} \cmidrule(lr){11-13} 
			& AME & SHAP & PDP & AME & SHAP & PDP & AME & SHAP & PDP & AME & SHAP & PDP \\
			\midrule
			age & 0.0002 & -0.0003 & 0.0002 & -0.0000 & 0.0011 & -0.0000 & -0.0004 & 0.0004 & -0.0002 & 0.0007 & -0.0010 & 0.0012 \\
			female & -0.0198 & -0.0005 & -0.0198 & -0.0062 & 0.0001 & -0.0062 & -0.0146 & 0.0002 & -0.0146 & -0.0135 & 0.0009 & -0.0200 \\
			education & 0.0002 & 0.0002 & 0.0002 & 0.0029 & 0.0006 & -0.0057 & 0.0030 & -0.0005 & -0.0356 & 0.0078 & -0.0016 & 0.0082 \\
			married & 0.0220 & -0.0038 & 0.0220 & 0.0054 & 0.0002 & 0.0054 & 0.0142 & 0.0002 & 0.0142 & 0.0192 & 0.0022 & 0.0217 \\
			\midrule
			avg\_pay\_status & -0.0256 & -0.0144 & -0.0246 & -0.0000 & 0.0029 & 0.0610 & 0.1027 & 0.0027 & 0.0404 & 0.0270 & -0.0116 & 0.0097 \\
			months\_delayed & 0.0826 & -0.0148 & 0.1107 & 0.0179 & -0.0008 & 0.0456 & 0.0209 & 0.0001 & 0.0404 & 0.1118 & -0.0154 & 0.0653 \\
			\midrule
			avg\_bill\_amt & 0.0000 & -0.0001 & 0.0000 & -0.0000 & -0.0018 & -0.0000 & -0.0000 & -0.0046 & -0.0000 & 0.0000 & 0.0036 & 0.0000 \\
			avg\_pay\_amt & -0.0000 & 0.0026 & -0.0000 & -0.0000 & -0.0006 & -0.0000 & -0.0000 & -0.0005 & -0.0000 & -0.0000 & -0.0025 & -0.0000 \\
			pay\_ratio & -0.0454 & 0.0100 & -0.0444 & -0.4306 & 0.0012 & -0.0998 & -0.4396 & 0.0005 & -0.0953 & 0.0009 & 0.0015 & -0.0170 \\
			\bottomrule
		\end{tabular}
		\begin{tablenotes}
			\footnotesize
			\item \textit{Notes:} AME = Average Marginal Effect (marginfx). SHAP = mean signed SHAP value (GradientExplainer). PDP = slope of partial dependence plot. Full specification (A+B+C) only. No standard errors reported for SHAP and PDP.
		\end{tablenotes}
	\end{threeparttable}
\end{sidewaystable}

Table~\ref{tab:credit_default_method_comparison} again reveals a systematic disconnect between SHAP values and AMEs. For the payment history variables the disagreement is particularly stark: the AME for months delayed is 0.083 for logistic regression while the mean signed SHAP value is $-$0.015 --- opposite in sign. The AME and PDP are in close agreement throughout, with the PDP slope for months delayed of 0.111 reasonably close to the AME of 0.083 given that PDP slopes are sensitive to the choice of evaluation grid. The pay ratio variable illustrates the most consequential divergence: the random forest and XGBoost AMEs are $-$0.431 and $-$0.440, indicating that a one-unit increase in the payment ratio is associated with a large reduction in default probability, while the corresponding SHAP values are 0.001 and 0.001 --- near zero and in the wrong direction. A researcher relying on SHAP alone would conclude that payment ratio is essentially irrelevant to default risk, which directly contradicts the AME and PDP evidence. And as before, with no standard errors attached to the SHAP or PDP estimates, there is no statistical basis for adjudicating between these conflicting signals.

\subsection{Ames Housing}

The third dataset is the Ames Housing dataset \citep{decock2011}, comprising approximately 1,460 observations of residential home sales in Ames, Iowa between 2006 and 2010. The outcome is sale price in dollars, so average marginal effects are interpreted in dollar units --- a one-unit change in a continuous feature is associated with a $\widehat{\mathrm{AME}}_j$ change in predicted sale price, on averege across the sample. Features are organized into physical characteristics (above-ground living area, number of bedrooms, number of bathrooms), quality and condition (overall quality rating, condition rating, year built, year remodeled), and lot and location (lot area, garage area, neighborhood tier indicators). The specification search considers linear regression and a random forest, covering the regression case and completing the sweep across outcome types.

\begin{table*}[htbp]
	\centering
	\small
	\resizebox{\textwidth}{!}{%
		\begin{threeparttable}
			\caption{Ames Housing: Specification Search. Average marginal effects on Sale Price (USD).}
			\label{tab:ames_housing_spec_search}
			\begin{tabular}{lrrrrrrrr}
				\toprule
				Variable & \multicolumn{4}{c}{Linear} & \multicolumn{4}{c}{Random Forest} \\
				\cmidrule(lr){2-5} \cmidrule(lr){6-9} 
				& Spec A & A+B & A+C & A+B+C & Spec A & A+B & A+C & A+B+C \\
				\midrule
				\multicolumn{9}{l}{\textit{Panel A: Physical Characteristics}} \\
				\quad GrLivArea & 108.22$^{***}$ & 85.99$^{***}$ & 74.32$^{***}$ & 65.57$^{***}$ & 126.41$^{***}$ & 79.37$^{***}$ & 78.92$^{***}$ & 56.44$^{***}$ \\
				& (13.45) & (15.05) & (11.03) & (13.84) & (8.22) & (3.47) & (3.61) & (1.93) \\
				\quad BedroomAbvGr & -27911.62$^{***}$ & -10181.55$^{***}$ & -11440.18$^{***}$ & -6189.02$^{**}$ & -23128.28$^{***}$ & -3055.72$^{***}$ & -5528.45$^{***}$ & -1470.18$^{***}$ \\
				& (3632.79) & (3107.43) & (3193.68) & (3042.41) & (1115.11) & (246.78) & (470.90) & (127.74) \\
				\quad FullBath & 30380.78$^{***}$ & -9793.90$^{*}$ & 5378.67 & -6287.42 & 18535.37$^{***}$ & 382.25$^{***}$ & 1687.49$^{***}$ & 498.18$^{***}$ \\
				& (5236.69) & (5750.79) & (3666.97) & (4593.14) & (601.00) & (104.69) & (108.04) & (73.44) \\
				\quad HalfBath & 3586.62 & -13111.01$^{***}$ & -728.29 & -6967.90$^{*}$ & 493.43 & -1160.44$^{***}$ & 420.55$^{***}$ & -198.68$^{***}$ \\
				& (4187.08) & (4093.96) & (2954.56) & (3607.14) & (491.19) & (111.48) & (97.27) & (51.07) \\
				\midrule
				\multicolumn{9}{l}{\textit{Panel B: Quality and Condition}} \\
				\quad OverallQual & -- & 21388.25$^{***}$ & -- & 16465.22$^{***}$ & -- & 17049.34$^{***}$ & -- & 15636.08$^{***}$ \\
				&  & (1824.85) &  & (1228.91) &  & (560.81) &  & (519.24) \\
				\quad OverallCond & -- & 5810.75$^{***}$ & -- & 6311.48$^{***}$ & -- & 2774.28$^{***}$ & -- & 1974.55$^{***}$ \\
				&  & (1112.27) &  & (934.11) &  & (126.97) &  & (89.75) \\
				\quad YearBuilt & -- & 685.31$^{***}$ & -- & 440.47$^{***}$ & -- & 729.64$^{***}$ & -- & 576.01$^{***}$ \\
				&  & (87.70) &  & (113.18) &  & (143.44) &  & (112.13) \\
				\quad YearRemodAdd & -- & 72.76 & -- & 101.51$^{*}$ & -- & 756.23$^{***}$ & -- & 518.52$^{***}$ \\
				&  & (62.62) &  & (53.22) &  & (72.12) &  & (46.58) \\
				\midrule
				\multicolumn{9}{l}{\textit{Panel C: Lot and Location}} \\
				\quad LotArea & -- & -- & 0.33 & 0.61$^{***}$ & -- & -- & 1.97$^{***}$ & 1.81$^{***}$ \\
				&  &  & (0.21) & (0.24) &  &  & (0.15) & (0.09) \\
				\quad GarageArea & -- & -- & 72.25$^{***}$ & 38.68$^{***}$ & -- & -- & 82.79$^{***}$ & 50.03$^{***}$ \\
				&  &  & (6.92) & (6.01) &  &  & (5.75) & (3.11) \\
				\quad nbhd\_high & -- & -- & 90316.66$^{***}$ & 51433.02$^{***}$ & -- & -- & 51108.94$^{***}$ & 5058.64$^{***}$ \\
				&  &  & (7451.69) & (9809.90) &  &  & (999.15) & (255.98) \\
				\quad nbhd\_mid & -- & -- & 31391.51$^{***}$ & 7508.07 & -- & -- & 25973.59$^{***}$ & 5949.97$^{***}$ \\
				&  &  & (3266.03) & (5151.49) &  &  & (531.38) & (259.64) \\
				\midrule
				$R^2$ & 0.584 & 0.758 & 0.739 & 0.802 & 0.829 & 0.937 & 0.926 & 0.953 \\
				RMSE & 51,243.9 & 39,080.1 & 40,596.4 & 35,355.5 & 32,830.9 & 19,881.6 & 21,534.2 & 17,244.4 \\
				\bottomrule
			\end{tabular}
			\begin{tablenotes}
				\footnotesize
				\item \textit{Notes:} $^{***}$p$<$0.01, $^{**}$p$<$0.05, $^{*}$p$<$0.10. Standard errors from nonparametric bootstrap (200 replicates) in parentheses.
			\end{tablenotes}
		\end{threeparttable}
	}%
\end{table*}

Table~\ref{tab:ames_housing_spec_search} presents the specification search. The regression setting allows AMEs to be interpreted directly in dollar units, which makes the substantive findings immediately legible. Above-ground living area carries an AME of \$108 per square foot in the demographics-only specification for linear regression, falling to \$66 in the full specification as quality, condition, and location controls absorb variation --- a pattern of positive omitted variable bias that makes intuitive sense, since larger homes tend to be in better neighborhoods and of higher quality. The random forest standard errors are dramatically tighter than those from linear regression throughout: the GrLivArea AME of \$56 in the full specification has a standard error of \$1.93 for the random forest versus \$13.84 for linear regression, a sevenfold reduction that reflects both the better model fit --- $R^2$ of 0.953 versus 0.802 --- and the smoother prediction surface of the ensemble. The overall quality rating is the largest single predictor in the quality panel, with an AME of approximately \$21,000 per rating point for linear regression and \$17,000 for the random forest in specification A+B, both precisely estimated. The neighborhood tier indicators show substantial attenuation as quality controls are added --- the high neighborhood AME falls from \$90,000 to \$51,000 for linear regression and from \$51,000 to \$5,000 for the random forest --- suggesting that neighborhood effects in the black-box model are largely mediated by the quality and condition variables in a way the linear model cannot fully represent.

\begin{table*}[htbp]
	\centering
	\small
	\begin{threeparttable}
		\caption{Ames Housing: Model Comparison. Average marginal effects on Sale Price (USD), full specification (A+B+C).}
		\label{tab:ames_housing_model_comparison}
		\begin{tabular}{lrrrr}
			\toprule
			Variable & Linear & Random Forest & XGBoost & Neural Net \\
			\midrule
			GrLivArea & 65.57$^{***}$ & 56.44$^{***}$ & 61.91$^{***}$ & 66.71$^{***}$ \\
			& (13.84) & (1.93) & (4.00) & (7.14) \\
			BedroomAbvGr & -6189.02$^{**}$ & -1470.18$^{***}$ & -2921.99$^{***}$ & -8026.62$^{***}$ \\
			& (3042.41) & (127.74) & (651.83) & (6.75) \\
			FullBath & -6287.42 & 498.18$^{***}$ & 372.18 & -874.83$^{***}$ \\
			& (4593.14) & (73.44) & (789.08) & (6.10) \\
			HalfBath & -6967.90$^{*}$ & -198.68$^{***}$ & -756.76$^{*}$ & -1967.74$^{***}$ \\
			& (3607.14) & (51.07) & (404.59) & (8.20) \\
			\midrule
			OverallQual & 16465.22$^{***}$ & 15636.08$^{***}$ & 13735.72$^{***}$ & 11140.42$^{***}$ \\
			& (1228.91) & (519.24) & (500.14) & (7.96) \\
			OverallCond & 6311.48$^{***}$ & 1974.55$^{***}$ & 4596.29$^{***}$ & 6074.62$^{***}$ \\
			& (934.11) & (89.75) & (533.35) & (6.22) \\
			YearBuilt & 440.47$^{***}$ & 576.01$^{***}$ & 619.02$^{***}$ & 335.19$^{***}$ \\
			& (113.18) & (112.13) & (125.80) & (3.21) \\
			YearRemodAdd & 101.51$^{*}$ & 518.52$^{***}$ & 647.81$^{***}$ & 612.99$^{***}$ \\
			& (53.22) & (46.58) & (187.76) & (4.04) \\
			\midrule
			LotArea & 0.61$^{***}$ & 1.81$^{***}$ & 2.11$^{***}$ & 2.28$^{***}$ \\
			& (0.24) & (0.09) & (0.27) & (0.53) \\
			GarageArea & 38.68$^{***}$ & 50.03$^{***}$ & 43.09$^{***}$ & 36.41$^{***}$ \\
			& (6.01) & (3.11) & (9.21) & (8.34) \\
			nbhd\_high & 51433.02$^{***}$ & 5058.64$^{***}$ & 20307.33$^{***}$ & 19166.75$^{***}$ \\
			& (9809.90) & (255.98) & (1644.21) & (17.29) \\
			nbhd\_mid & 7508.07 & 5949.97$^{***}$ & 9744.98$^{***}$ & 12711.08$^{***}$ \\
			& (5151.49) & (259.64) & (557.43) & (9.87) \\
			\midrule
			$R^2$ & 0.802 & 0.953 & 0.954 & 0.939 \\
			RMSE & 35,355.5 & 17,244.4 & 17,036.6 & 19,574.7 \\
			\bottomrule
		\end{tabular}
		\begin{tablenotes}
			\footnotesize
			\item \textit{Notes:} $^{***}$p$<$0.01, $^{**}$p$<$0.05, $^{*}$p$<$0.10. Standard errors from nonparametric bootstrap (200 replicates) in parentheses.
		\end{tablenotes}
	\end{threeparttable}
\end{table*}

Table~\ref{tab:ames_housing_model_comparison} compares all four models on the full specification. All four models are well-fitted, with $R^2$ ranging from 0.802 for linear regression to 0.954 for XGBoost and the neural network reaching 0.939 with normalization. The core variables show striking agreement across architectures. The GrLivArea AME is \$66, \$56, \$62, and \$67 for linear regression, random forest, XGBoost, and neural network respectively --- a range of \$11 across four models that differ fundamentally in their functional form. The overall quality AME ranges from \$11,140 to \$16,465, all precisely estimated and within the same order of magnitude. This stability is not guaranteed by the method --- it reflects genuine agreement across models about the underlying conditional expectation function --- and it is the kind of cross-model validation that the AME framework makes possible precisely because it produces output on a common scale. The bedroom variable illustrates genuine model disagreement: the AME is $-$\$6,189 for linear regression, $-$\$1,470 for random forest, $-$\$2,922 for XGBoost, and $-$\$8,027 for the neural network, all negative and significant. The negative bedroom effect, conditional on living area, is a well-known finding in the hedonic pricing literature --- holding square footage fixed, more bedrooms means smaller rooms --- and all four models recover it, though with magnitudes spanning a fivefold range. The neural network standard errors are implausibly tight throughout, on the order of single dollars for several variables, which is a known artifact of deterministic gradient computation through a fixed trained network rather than a genuine reflection of estimation uncertainty; the bootstrap resampling propagates uncertainty through the model fitting but the neural network's smooth differentiable surface produces nearly identical finite differences across replicates.

\begin{sidewaystable}[htbp]
	\centering
	\small
	\begin{threeparttable}
		\caption{Ames Housing: Method Comparison. AME, SHAP, and PDP estimates for Sale Price (USD), full specification (A+B+C).}
		\label{tab:ames_housing_method_comparison}
		\begin{tabular}{lrrrrrrrrrrrr}
			\toprule
			Variable & \multicolumn{3}{c}{Linear} & \multicolumn{3}{c}{Random Forest} & \multicolumn{3}{c}{XGBoost} & \multicolumn{3}{c}{Neural Net} \\
			\cmidrule(lr){2-4} \cmidrule(lr){5-7} \cmidrule(lr){8-10} \cmidrule(lr){11-13} 
			& AME & SHAP & PDP & AME & SHAP & PDP & AME & SHAP & PDP & AME & SHAP & PDP \\
			\midrule
			GrLivArea & 65.57 & 3423.49 & 65.57 & 56.44 & 61.69 & 38.64 & 61.91 & -971.88 & 39.06 & 66.71 & -- & 55.83 \\
			BedroomAbvGr & -6189.02 & 145.82 & -6189.02 & -1470.18 & 90.87 & -1272.27 & -2921.99 & 158.44 & -4456.61 & -8026.62 & -- & -7803.94 \\
			FullBath & -6287.42 & 31.01 & -6287.42 & 498.18 & 34.06 & 543.63 & 372.18 & 54.86 & 193.04 & -874.83 & -- & 339.37 \\
			HalfBath & -6967.90 & -298.76 & -6967.90 & -198.68 & 50.46 & -125.75 & -756.76 & 30.53 & -1495.20 & -1967.74 & -- & -3985.01 \\
			\midrule
			OverallQual & 16465.22 & 1635.24 & 16465.22 & 15636.08 & 596.70 & 20776.14 & 13735.72 & 713.13 & 20130.92 & 11140.42 & -- & 11855.02 \\
			OverallCond & 6311.48 & 980.44 & 6311.48 & 1974.55 & 161.26 & 2030.44 & 4596.29 & -481.60 & 5071.47 & 6074.62 & -- & 6504.98 \\
			YearBuilt & 440.47 & -415.01 & 440.47 & 576.01 & 288.29 & 209.05 & 619.02 & 569.00 & 397.33 & 335.19 & -- & 415.03 \\
			YearRemodAdd & 101.51 & 86.86 & 101.51 & 518.52 & 373.04 & 118.06 & 647.81 & 499.84 & 152.00 & 612.99 & -- & 475.07 \\
			\midrule
			LotArea & 0.61 & 267.44 & 0.61 & 1.81 & -183.50 & 1.33 & 2.11 & 216.28 & 2.38 & 2.28 & -- & 2.24 \\
			GarageArea & 38.68 & 93.22 & 38.68 & 50.03 & -884.31 & 23.77 & 43.09 & -704.59 & 28.98 & 36.41 & -- & 47.34 \\
			nbhd\_high & 51433.02 & 77.50 & 51433.02 & 5058.64 & -506.36 & 5058.64 & 20307.33 & 50.70 & 20307.34 & 19166.75 & -- & 6383.02 \\
			nbhd\_mid & 7508.07 & 266.38 & 7508.07 & 5949.97 & 40.68 & 5949.97 & 9744.98 & -134.71 & 9744.97 & 12711.08 & -- & 9748.84 \\
			\bottomrule
		\end{tabular}
		\begin{tablenotes}
			\footnotesize
			\item \textit{Notes:} AME = Average Marginal Effect (marginfx). SHAP = mean signed SHAP value (GradientExplainer). PDP = slope of partial dependence plot. Full specification (A+B+C) only. No standard errors reported for SHAP and PDP.
		\end{tablenotes}
	\end{threeparttable}
\end{sidewaystable}

Table~\ref{tab:ames_housing_method_comparison} brings the SHAP comparison into sharpest relief in the regression setting, where the units are dollars and the magnitudes are large enough to make the divergences concrete. SHAP values are unavailable for the neural network in this table, reflecting a limitation of the GradientExplainer implementation for this architecture. For the other three models the divergence between SHAP and AME is stark. The SHAP value for GrLivArea is \$3,423 for the linear model, which attributes to that variable an effect fifty times larger than the AME of \$66 --- a number that is not interpretable as a marginal effect in any standard sense. The SHAP values for garage area ($-$\$884), lot area ($-$\$184), and overall condition ($-$\$482) are negative for the random forest while the AME and PDP estimates are positive for all three. The AME and PDP estimates are in close agreement throughout for all four models --- the neighborhood tier indicators match almost exactly between AME and PDP for random forest and XGBoost, the bedroom AME of $-$\$8,027 for the neural network is close to its PDP slope of $-$\$7,804, and the quality variables align across estimators within reasonable range. The Ames Housing results complete the sweep across outcome types: the framework produces dollar-denominated marginal effects with valid standard errors for a continuous regression outcome, using exactly the same estimation procedure as the binary classification examples, with no modification required.

\section{Simulation}
\label{sec:simulation}

We evaluate finite-sample performance through two simulation studies. The full
design and complete results are reported in the Supplementary Material
\citep{parker2026supp}; this
section summarises the findings.

Simulation~1 examines bias and root mean squared error of the AME estimates as a
function of sample size. Four features are drawn independently from $N(0,1)$,
with $x_1$ and $x_2$ carrying nonzero true effects and $x_3, x_4$ pure noise.
Three data generating processes --- linear, nonlinear ($\eta = 2x_1^2 + 3x_2$),
and interaction ($\eta = 2x_1 + 3x_2 + 2x_1 x_2$) --- are crossed with binary and
continuous outcomes, and four model classes are fit: a parametric baseline
(linear or logistic regression), a feedforward neural network, XGBoost, and a
random forest. Sample sizes range over $n \in \{250, \ldots, 5\,000\}$ with 500
to 1{,}000 Monte Carlo iterations per cell, and true AMEs are obtained by Monte
Carlo integration at $n = 10^6$.

Three findings emerge. First, the AME estimator is consistent for the parametric
baselines, neural networks, and XGBoost across all three data generating
processes and both outcome types, with RMSE declining at the expected
$O(n^{-1/2})$ rate. Second, random forests exhibit persistent negative bias that
does not vanish with sample size, reaching $-0.185$ in the nonlinear regression
design at $n = 5\,000$; this reflects the shrinkage imposed by depth and
subsampling constraints at default tuning. Third, AMEs are robust to interaction
misspecification when the interacting features have zero mean --- a linear model
omitting the $2x_1 x_2$ term still recovers the correct marginal effects, since
the interaction contributes $\mathbb{E}[2x_2] = 0$ to the marginal effect of
$x_1$.

Simulation~2 evaluates bootstrap calibration by comparing empirical coverage of
nominal 95\% confidence intervals against the target, using 200 bootstrap
replicates per Monte Carlo iteration. Coverage is close to nominal for the
parametric baselines, neural networks, and XGBoost, with departures at small
samples falling in the conservative direction. Random forests are again the
exception: coverage collapses at $n = 5\,000$, falling to 0.130 for $x_1$ in the
classification design and 0.104 for $x_1$ in the regression design. The mechanism
is the bias documented in Simulation~1. The bootstrap intervals contract at the
correct rate but around an off-center point estimate, so by $n = 5\,000$ the
interval width is comparable to the magnitude of the bias itself and coverage of
the true AME becomes nearly impossible. This is the empirical counterpart to the
theoretical status of bootstrap validity for Breiman forests established in
Section~\ref{sec:theory}, where the result is stated as a conjecture rather than
a theorem.

\section{Conclusion}
\label{sec:conclusion}

This paper has proposed a framework for interpreting any supervised machine learning
model through the lens of average marginal effects. The central claim is that the
traditional accuracy-interpretability tradeoff is not fundamental. By defining the
average marginal effect as the expectation of the partial derivative of the prediction
function with respect to each feature, averaged over the empirical distribution of
the covariates, we obtain an estimand that is well-defined for any supervised learner
and reduces to the classical econometric AME when the model is a generalized linear
model. The estimator replaces analytical derivatives with adaptive centered finite
differences, applies uniformly across model classes and outcome types, and produces a
coefficient table --- estimate, standard error, and significance level --- directly
analogous to the output of a parametric regression. The nonparametric bootstrap, with
full model refitting and hyperparameter selection on each replicate, provides standard
errors that capture uncertainty about model specification as well as estimation. These
are capabilities that existing interpretability methods, including SHAP
\citep{lundberg2017shap}, LIME \citep{ribeiro2016lime}, and partial dependence plots
\citep{friedman2001pdp}, do not provide.

The theoretical contribution proceeds in two parts. The first establishes gradient
consistency for three classes of supervised learners. For neural networks, the
appropriate function space is the Sobolev space $W^{1,p}$, and gradient consistency
follows from uniform $W^{2,p}$ boundedness combined with $L^p$ consistency via the
Rellich--Kondrachov compact embedding theorem \citep{evans2010pde}. For gradient
boosted models with tree base learners and for Mondrian forests
\citep{mourtada2017mondrian}, the appropriate function space is the space of functions
of bounded variation \citep{ambrosio2000bv}, where gradients exist as Radon measures
and consistency is established via the BV compactness theorem. For Breiman random
forests \citep{breiman2001rf, scornet2015forests}, gradient consistency is conjectured
rather than proved: the obstruction lies in the data-dependent partition boundaries,
whose discontinuity structure does not admit the geometric analysis available for
Mondrian partitions. The second part establishes bootstrap validity for all three
proved estimator classes --- neural networks, gradient boosted models, and Mondrian
forests --- via the Hadamard differentiability of the AME functional
\citep{vandervaart1996weak} and the functional delta method. Corollaries~1 through~3
make this precise, showing that the bootstrap distribution of the sample AME
consistently estimates the sampling distribution of the estimator under rate
conditions that are verified for each class. Bootstrap validity for Breiman forests
remains open for the same reason gradient consistency is harder to establish ---
the data-dependent partition geometry resists the clean decomposition that the
Mondrian independence assumption makes available.

The empirical results across three datasets --- UCI Adult Income \citep{kohavi1996,
	UCI}, UCI Credit Card Default \citep{yeh2009, UCI}, and Ames Housing
\citep{decock2011} --- validate the framework on binary classification and continuous
regression outcomes. The AME estimates are stable across model classes where the
underlying models are well-fitted, the bootstrap standard errors narrow as model fit
improves, and the AME and partial dependence plot estimates are in close agreement
throughout. The comparison with SHAP values reveals a systematic divergence: SHAP
values are on a different scale, in several cases have the wrong sign relative to AME
and PDP, and carry no inferential apparatus. The Ames Housing application is
particularly informative because the dollar-denominated units make the divergence
concrete --- SHAP attributes to above-ground living area an effect sixty times larger
than the AME for the linear model, a number that is not interpretable as a marginal
effect in any standard sense. An additional finding is that Breiman random forests
produce less stable AME estimates than the other estimator classes, consistent with
the theoretical result that gradient consistency for Breiman forests is harder to
establish. Practitioners requiring random forest AMEs are advised to consider Mondrian
forests, for which both gradient consistency and bootstrap validity are proved.

Several important problems remain open. The most pressing is gradient consistency and
bootstrap validity for Breiman forests. The obstruction identified in the proof of
Theorem~3 --- that data-dependent splits couple partition geometry to response
variation in a way that prevents the key decomposition --- suggests that new
analytical tools will be required, possibly involving direct control over the
data-dependent splitting criterion or a coupling argument relating Breiman partitions
to Mondrian partitions in an appropriate sense. Establishing these results would
close the gap between the theoretical guarantees and the most widely used random
forest implementation.

A second open problem is the extension of gradient consistency and bootstrap validity
to additional model classes. The theoretical results cover neural networks, gradient
boosted models, and Mondrian forests. The BV and Sobolev architectures developed here
provide a template: identify the natural function space for the estimator class,
establish derivative existence and consistency in that space, verify the rate
condition, and invoke the Hadamard differentiability of the AME functional. Applying
this template to support vector machines with smooth kernels, Bayesian additive
regression trees, generalized additive models, and kernel methods is an important
direction. Each case requires identifying the appropriate function space and verifying
the well-behaved conditions, but the framework makes the proof obligations explicit
and the path tractable. Splines and ridge regression are particularly natural
candidates since they live in Sobolev spaces by construction and the gradient
consistency results are likely near-immediate from the existing theory.

A third open problem is the extension to causal inference. The average marginal effect
as defined here is a predictive quantity --- the average derivative of the conditional
expectation function with respect to $x_j$, under the observational distribution of
the covariates. It is not in general a causal effect. Extending the framework to
settings where causal identification is possible, for example through instrumental
variables or regression discontinuity designs, would substantially broaden its
applicability to policy-relevant questions. The honest forest framework of
\citet{wager2018forests} and the broader literature on causal machine learning suggest
a natural path forward, but the connection to the AME estimand studied here has not
been formalized.

A fourth direction concerns the discrete and mixed-type feature setting. The current
framework handles binary variables through a finite difference contrast and continuous
variables through an adaptive centered finite difference, but ordered categorical
variables, count variables, and interactions among features of mixed type are not
treated systematically. The marginal effects at the mean versus average marginal
effects distinction \citep{greene2012, wooldridge2010} also warrants further study in
the nonparametric setting, where the mean of the covariates may lie in a sparse region
of the feature space and the marginal effect at the mean may therefore be a poor
summary of the average effect.

Finally, the computational cost of the full bootstrap procedure --- refitting the
model with hyperparameter selection on each of $B$ replicates --- limits practical
applicability for very large datasets or very expensive models. Approximations based
on the infinitesimal jackknife \citep{wager2014confidence}, subsampling, or
warm-started hyperparameter search reduce this cost substantially, but the statistical
properties of these approximations relative to the full bootstrap have not been
characterized. The Hadamard differentiability established in Lemma~5 provides a
natural starting point for analyzing such approximations: any procedure that
consistently estimates the derivative $T'_f$ of the AME functional will yield
asymptotically valid inference, and this opens a principled path toward
computationally efficient alternatives to the full bootstrap. Developing and
characterizing such procedures is an important practical problem that we leave to
future work.

\begin{acks}[Acknowledgments]
The author is grateful to Donggyu Sul for extensive comments and discussion
that substantially improved this work. Helpful comments were also provided by
Walter Johnston, Max Pariazar and Rahul Sethi. Any remaining errors are the
author's own.
\end{acks}

\begin{supplement}
\stitle{Supplement to ``Model-Agnostic Interpretation of Machine Learning via
Average Marginal Effects'': Proofs and full simulation results}
\sdescription{Contains complete proofs of all lemmas, theorems and corollaries
stated in Section~\ref{sec:theory} (Appendices~A--L), together with the full
design and complete results of the simulation studies summarised in
Section~\ref{sec:simulation} (Appendix~M).}
\slink[doi]{COMPLETED BY TYPESETTER}
\sdatatype{.pdf}
\end{supplement}

\bibliographystyle{imsart-nameyear}
\bibliography{references}

\end{document}
